\documentclass[11pt]{article}

% ==================== PAQUETES ====================
\usepackage[utf8]{inputenc}
\usepackage[T1]{fontenc}
\IfFileExists{spanish.ldf}{\usepackage[spanish,es-tabla]{babel}}{}
\usepackage{amsmath,amssymb,amsthm,mathtools}
\usepackage{bm}
\usepackage[margin=2.5cm]{geometry}
\usepackage{enumitem}
\usepackage{xcolor}
\usepackage{hyperref}
\hypersetup{colorlinks=true,linkcolor=black,urlcolor=black,citecolor=black}

% ==================== NOTACIÓN ====================
\newcommand{\N}{\mathcal{N}}
\newcommand{\IG}{\mathcal{IG}}
\newcommand{\Ga}{\mathcal{G}}
\newcommand{\R}{\mathbb{R}}
\newcommand{\E}{\mathbb{E}}
\newcommand{\Var}{\operatorname{Var}}
\newcommand{\Cov}{\operatorname{Cov}}
\newcommand{\tr}{\operatorname{tr}}
\newcommand{\vecop}{\operatorname{vec}}
\newcommand{\diag}{\operatorname{diag}}
\newcommand{\ind}{\mathbb{1}}
\newcommand{\bpar}{\bm{\beta}}
\newcommand{\bX}{\mathbf{X}}
\newcommand{\by}{\mathbf{y}}
\newcommand{\bY}{\mathbf{Y}}
\newcommand{\bz}{\mathbf{z}}
\newcommand{\bB}{\mathbf{B}}
\newcommand{\bI}{\mathbf{I}}
\newcommand{\bV}{\mathbf{V}}
\newcommand{\bW}{\mathbf{W}}
\newcommand{\bSig}{\bm{\Sigma}}
\newcommand{\bmu}{\bm{\mu}}
\newcommand{\beps}{\bm{\varepsilon}}
\newcommand{\bOm}{\bm{\Omega}}
\newcommand{\kron}{\otimes}
\newcommand{\dpar}{\mathrm{d}}
\DeclareMathOperator*{\argmax}{arg\,max}

\theoremstyle{definition}
\newtheorem{prop}{Proposición}

\title{\textbf{Inferencia Bayesiana}\\[2pt]\large Notas de clase\footnote{Basado en mis notas de clase de los cursos \emph{Inferencia Bayesiana} y \emph{Econometría Macro Avanzada} de la maestría de la PUCP.}}
\author{Josué David Cárdenas Barrantes}
\date{}

\begin{document}
\maketitle
\tableofcontents
\newpage

% ============================================================
\section{Repaso de probabilidades}
% ============================================================
Una probabilidad es una transformación $P$ que asigna a cada evento $A$ de un espacio muestral $\Omega$ un número real $P(A)$ y que satisface los tres axiomas de la probabilidad:
\begin{itemize}[leftmargin=2em]
  \item[\textbf{A1}] Para cualquier evento $A$: $P(A)\ge 0$.
  \item[\textbf{A2}] $P(\Omega)=1$.
  \item[\textbf{A3}] Para eventos disjuntos $A_1,A_2,\dots$: $P\!\left(\bigcup_i A_i\right)=\sum_i P(A_i)$.
\end{itemize}

% ============================================================
\section{Idea general}
% ============================================================
En el enfoque bayesiano, la pregunta que nos hacemos es diferente al enfoque frecuentista: \emph{¿qué tan probable es cada posible valor de $\beta$, dado que observé los datos?} La respuesta es una distribución completa.
\[
p(\beta\mid y)=\frac{p(y\mid\beta)\,p(\beta)}{p(y)}
\]
\begin{itemize}[leftmargin=2em]
  \item $p(\beta\mid y)=$ lo que creo sobre $\beta$ después de ver los datos (\textbf{posterior}).
  \item $p(y\mid\beta)=$ \textbf{verosimilitud}: qué tan probables son los datos dado el parámetro $\beta$.
  \item $p(\beta)=$ \textbf{prior}: lo que creo sobre $\beta$ antes de ver los datos.
  \item $p(y)=$ constante normalizadora.
\end{itemize}
\[
p(\beta\mid y)\ \propto\ p(y\mid\beta)\,p(\beta)
\]

% ============================================================
\section{Regresión lineal bayesiana}
% ============================================================
Se define el modelo de regresión simple
\[
y_t = X_t'\beta+\varepsilon_t,
\]
donde $\varepsilon_1,\dots,\varepsilon_T$ son errores i.i.d.\ $(0,\sigma^2)$, $X_t=(1,X_{1,t},\dots,X_{K-1,t})'$ son $K$ regresores y $\beta=(\beta_0,\beta_1,\dots,\beta_{K-1})'$ son los estimadores asociados.

\subsection{Función de verosimilitud}
En forma matricial,
\[
\underbrace{\begin{bmatrix}y_1\\ y_2\\ \vdots\\ y_T\end{bmatrix}}_{T\times 1}
=
\underbrace{\begin{bmatrix}
1 & X_{1,1} & \cdots & X_{K-1,1}\\
1 & X_{1,2} & \cdots & X_{K-1,2}\\
\vdots & \vdots & & \vdots\\
1 & X_{1,T} & \cdots & X_{K-1,T}
\end{bmatrix}}_{T\times K}
\underbrace{\begin{bmatrix}\beta_0\\ \beta_1\\ \vdots\\ \beta_{K-1}\end{bmatrix}}_{K\times1}
+
\underbrace{\begin{bmatrix}\varepsilon_1\\ \varepsilon_2\\ \vdots\\ \varepsilon_T\end{bmatrix}}_{T\times1},
\qquad
\bY_{T\times1}=\bX_{T\times K}\,\bpar_{K\times1}+\beps_{T\times1}.
\]
Los errores tienen distribución normal multivariada $\beps\sim\N(\mathbf{0}_{T\times1},\sigma^2\bI_T)$. Si $X$ es exógeno, entonces $(\by\mid\beta,\sigma^2)\sim\N(\bX\beta,\sigma^2\bI_T)$, y la verosimilitud es
\begin{align*}
\mathcal{L}(\beta,\sigma^2)=p(\by\mid\beta,\sigma^2)
&=|2\pi\sigma^2\bI_T|^{-1/2}\exp\!\left(-\tfrac{1}{2\sigma^2}(\by-\bX\beta)'(\by-\bX\beta)\right)\\
&=(2\pi\sigma^2)^{-T/2}\exp\!\left(-\tfrac{1}{2\sigma^2}(\by-\bX\beta)'(\by-\bX\beta)\right)\\
&\propto (\sigma^2)^{-T/2}\exp\!\left(-\tfrac{1}{2\sigma^2}(\by-\bX\beta)'(\by-\bX\beta)\right),
\end{align*}
donde $|\cdot|$ denota el determinante y se cumple que, para un escalar $c$ y una matriz $A$ de dimensión $T\times T$, $|cA|=c^{T}|A|$.

\subsection{Descomposición OLS}
Los estimadores OLS de $\beta$ y $\sigma^2$ son
\[
\hat\beta=(\bX'\bX)^{-1}(\bX'\by),\qquad
s^2=\frac{(\by-\bX\hat\beta)'(\by-\bX\hat\beta)}{T-K}\ \ (\text{ajustado por grados de libertad}).
\]
Expandiendo el término cuadrático,
\begin{align*}
(\by-\bX\beta)'(\by-\bX\beta)
&=(\by-\bX\hat\beta+\bX\hat\beta-\bX\beta)'(\by-\bX\hat\beta+\bX\hat\beta-\bX\beta)\\
&=\big[(\by-\bX\hat\beta)-\bX(\beta-\hat\beta)\big]'\big[(\by-\bX\hat\beta)-\bX(\beta-\hat\beta)\big]\\
&=(\by-\bX\hat\beta)'(\by-\bX\hat\beta)+(\beta-\hat\beta)'\bX'\bX(\beta-\hat\beta)\\
&=(T-K)\,s^2+(\beta-\hat\beta)'\bX'\bX(\beta-\hat\beta),
\end{align*}
donde los términos cruzados se anulan porque
\[
\bX'(\by-\bX\hat\beta)=\bX'\big(\by-\bX(\bX'\bX)^{-1}\bX'\by\big)=\big(\bX'-\bX'\bX(\bX'\bX)^{-1}\bX'\big)\by=\mathbf{0}.
\]
Reemplazando en la verosimilitud,
\[
\mathcal{L}(\beta,\sigma^2)\propto(\sigma^2)^{-T/2}\exp\!\left(-\tfrac12\Big[(T-K)\tfrac{s^2}{\sigma^2}+(\beta-\hat\beta)'\tfrac{\bX'\bX}{\sigma^2}(\beta-\hat\beta)\Big]\right).
\]

\subsection{Distribuciones prior}
Se asumen priors independientes $\beta\sim\N(\underline{\beta},\underline{\bV}_\beta)$ ($K\times1$, $K\times K$) y $\sigma^2\sim\IG(\underline{\nu},\underline{S})$:
\begin{align*}
p(\beta)&=(2\pi)^{-K/2}|\underline{\bV}_\beta|^{-1/2}\exp\!\left(-\tfrac12(\beta-\underline{\beta})'\underline{\bV}_\beta^{-1}(\beta-\underline{\beta})\right)
\propto\exp\!\left(-\tfrac12(\beta-\underline{\beta})'\underline{\bV}_\beta^{-1}(\beta-\underline{\beta})\right),\\[4pt]
p(\sigma^2)&=\frac{\underline{S}^{\underline{\nu}}}{\Gamma(\underline{\nu})}(\sigma^2)^{-(\underline{\nu}+1)}\exp\!\left(-\tfrac{1}{\sigma^2}\underline{S}\right)
\propto(\sigma^2)^{-(\underline{\nu}+1)}\exp\!\left(-\tfrac{1}{\sigma^2}\underline{S}\right).
\end{align*}
La constante de $p(\beta)$ se obvia porque no depende de $\beta$; queda el kernel de una normal. Análogamente $p(\sigma^2)$ es el kernel de una Inverse-Gamma. La prior conjunta (por independencia) es
\[
p(\beta,\sigma^2)=p(\beta)\,p(\sigma^2)
\propto\Big\{\exp\!\big(-\tfrac12(\beta-\underline{\beta})'\underline{\bV}_\beta^{-1}(\beta-\underline{\beta})\big)\Big\}
\Big\{(\sigma^2)^{-(\underline{\nu}+1)}\exp\!\big(-\tfrac{1}{\sigma^2}\underline{S}\big)\Big\}.
\]

\subsection{Distribución posterior conjunta}
\begin{align*}
p(\beta,\sigma^2\mid\by)&\propto p(\by\mid\beta,\sigma^2)\,p(\beta,\sigma^2)\\
&\propto\Big\{(\sigma^2)^{-T/2}\exp\!\big(-\tfrac12[(T-K)\tfrac{s^2}{\sigma^2}+(\beta-\hat\beta)'\tfrac{\bX'\bX}{\sigma^2}(\beta-\hat\beta)]\big)\Big\}\\
&\quad\times\Big\{\exp\!\big(-\tfrac12(\beta-\underline{\beta})'\underline{\bV}_\beta^{-1}(\beta-\underline{\beta})\big)\times(\sigma^2)^{-(\underline{\nu}+1)}\exp\!\big(-\tfrac{1}{\sigma^2}\underline{S}\big)\Big\}.
\end{align*}
Esta expresión no tiene forma cerrada, por lo que se simula mediante los algoritmos de \textbf{Gibbs Sampling} y \textbf{Metropolis--Hastings}.

\subsection{Distribuciones condicionales (Gibbs Sampling)}

\paragraph{Condicional de $\sigma^2$ (fijando $\beta$).}
\begin{align*}
p(\sigma^2\mid\beta,\by)&=p(\by\mid\beta,\sigma^2)\,p(\sigma^2\mid\beta)\\
&\propto(\sigma^2)^{-T/2}\exp\!\big(-\tfrac{1}{2\sigma^2}(\by-\bX\beta)'(\by-\bX\beta)\big)\times(\sigma^2)^{-(\underline{\nu}+1)}\exp\!\big(-\tfrac{1}{\sigma^2}\underline{S}\big)\\
&\propto(\sigma^2)^{-(\frac{T}{2}+\underline{\nu}+1)}\exp\!\left(-\tfrac{1}{\sigma^2}\Big[\underline{S}+\tfrac12(\by-\bX\beta)'(\by-\bX\beta)\Big]\right).
\end{align*}
La forma funcional es la de una Inverse-Gamma:
\[
p(\sigma^2\mid\beta,\by)\sim\IG\!\left(\tfrac{T}{2}+\underline{\nu},\ \underline{S}+\tfrac12(\by-\bX\beta)'(\by-\bX\beta)\right)\equiv\IG(\bar\nu,\bar S),
\]
con $\bar\nu=\underline{\nu}+\tfrac{T}{2}$ y $\bar S=\underline{S}+\tfrac12(\by-\bX\beta)'(\by-\bX\beta)$.

\paragraph{Condicional de $\beta$ (fijando $\sigma^2$).}
\begin{align*}
p(\beta\mid\sigma^2,\by)
&\propto\exp\!\left(-\tfrac12\Big[(\beta-\hat\beta)'\tfrac{\bX'\bX}{\sigma^2}(\beta-\hat\beta)+(\beta-\underline{\beta})'\underline{\bV}_\beta^{-1}(\beta-\underline{\beta})\Big]\right).
\end{align*}
Expandiendo los términos cuadráticos,
\begin{align*}
(\beta-\underline{\beta})'\underline{\bV}_\beta^{-1}(\beta-\underline{\beta})&=\beta'\underline{\bV}_\beta^{-1}\beta-2\beta'\underline{\bV}_\beta^{-1}\underline{\beta}+\text{cte},\\
(\beta-\hat\beta)'\tfrac{\bX'\bX}{\sigma^2}(\beta-\hat\beta)&=\beta'\tfrac{\bX'\bX}{\sigma^2}\beta-2\beta'\tfrac{\bX'\bX}{\sigma^2}\hat\beta+\text{cte},
\end{align*}
y reemplazando,
\begin{align*}
p(\beta\mid\sigma^2,\by)
&\propto\exp\!\left(-\tfrac12\Big[\beta'\underline{\bV}_\beta^{-1}\beta-2\beta'\underline{\bV}_\beta^{-1}\underline{\beta}+\beta'\tfrac{\bX'\bX}{\sigma^2}\beta-2\beta'\tfrac{\bX'\bX}{\sigma^2}\hat\beta\Big]\right)\\
&\propto\exp\!\left(-\tfrac12\Big[\beta'\big(\underline{\bV}_\beta^{-1}+\tfrac{1}{\sigma^2}\bX'\bX\big)\beta-2\beta'\big(\underline{\bV}_\beta^{-1}\underline{\beta}+\tfrac{1}{\sigma^2}\bX'\by\big)\Big]\right).
\end{align*}

\paragraph{Identificación de la normal.}
Para una gaussiana multivariada $\beta\sim\N(\mu,\Sigma)$ el exponente es
$(\beta-\mu)'\Sigma^{-1}(\beta-\mu)=\beta'\Sigma^{-1}\beta-2\beta'\Sigma^{-1}\mu+\text{cte}$. Comparando,
\[
\bar\bV_\beta=\Big(\underline{\bV}_\beta^{-1}+\tfrac{1}{\sigma^2}\bX'\bX\Big)^{-1},\qquad
\bar\beta=\bar\bV_\beta\Big(\underline{\bV}_\beta^{-1}\underline{\beta}+\tfrac{1}{\sigma^2}\bX'\by\Big),
\]
de modo que
\[
p(\beta\mid\sigma^2,\by)\sim\N(\bar\beta,\bar\bV_\beta).
\]

\subsection{Algoritmo de Gibbs Sampling}
Proponemos un punto inicial $\beta^{(0)}=a_0$ y $\sigma^{2(0)}=b_0>0$ (condiciones iniciales). Sampleamos $i=N+B$ iteraciones:
\begin{enumerate}[leftmargin=2em]
  \item Simular $\beta^{(i)}$ de $p(\beta^{(i)}\mid\sigma^{2(i-1)},\by)\sim\N(\bar\beta,\bar\bV_\beta)$.
  \item Simular $\sigma^{2(i)}$ de $p(\sigma^{2(i)}\mid\beta^{(i)},\by)\sim\IG(\bar\nu,\bar S)$.
\end{enumerate}
Para evitar la influencia de las condiciones iniciales se \emph{queman} las primeras $B$ simulaciones (\emph{burn-in}) hasta llegar a la distribución posterior; se conservan las $N$ finales.

\subsection{Prior menos informativa}
En la medida que queramos que la prior $\beta\sim\N(\underline{\beta},\underline{\bV}_\beta)$ sea menos informativa, elevamos el valor de la varianza. A partir de
\[
\bar\beta=\Big[\underline{\bV}_\beta^{-1}+\tfrac{1}{\sigma^2}\bX'\bX\Big]^{-1}\Big[\underline{\bV}_\beta^{-1}\underline{\beta}+\tfrac{1}{\sigma^2}\bX'\by\Big],
\]
cuando $\underline{\bV}_\beta\to\infty$ se tiene $\underline{\bV}_\beta^{-1}\to 0$, por lo tanto
\[
\bar\beta=\Big[\tfrac{1}{\sigma^2}\bX'\bX\Big]^{-1}\Big[\tfrac{1}{\sigma^2}\bX'\by\Big]=(\bX'\bX)^{-1}\bX'\by=\hat\beta_{\text{OLS}}.
\]
Mientras más pequeña sea la varianza de la prior, más informativa es y más pesa la prior; mientras más grande, más peso tiene la verosimilitud. En el límite, la prior no aporta nada y el estimador coincide con OLS/ML.

% ============================================================
\section{Herramientas para BVAR}
% ============================================================
\subsection{Vectorización de matrices}
La operación $\vecop(A)$ toma una matriz $m\times n$ y la convierte en un vector columna $mn\times1$ apilando las columnas. Si $A=\begin{bmatrix}a_{11}&a_{12}\\ a_{21}&a_{22}\end{bmatrix}$, entonces $\vecop(A)=(a_{11},a_{21},a_{12},a_{22})'$.

\paragraph{Propiedad 1: Linealidad.}
$\vecop(A+B)=\vecop(A)+\vecop(B)$, \quad $\vecop(\alpha A)=\alpha\,\vecop(A)$.

\paragraph{Propiedad 2: Puente con Kronecker.}
Para matrices $A,B,C$:
\[
\vecop(ABC)=(C^{\top}\kron A)\,\vecop(B).
\]
Casos especiales: si $C=I$, $\vecop(AB)=(I\kron A)\vecop(B)$; si $A=I$, $\vecop(BC)=(C^{\top}\kron I)\vecop(B)$.

\paragraph{Propiedad 3: Traza y formas cuadráticas.}
\[
\tr(A^{\top}B)=\vecop(A)^{\top}\vecop(B),\qquad
\vecop(B)^{\top}(C\kron A)\vecop(B)=\tr(B^{\top}ABC).
\]

\paragraph{Producto de Kronecker.}
$A_{m\times n}\kron B_{p\times q}=(A\kron B)_{mp\times nq}$: las dimensiones se multiplican. Por ejemplo $I_n\kron X_{T\times K}=(I_n\kron X)_{nT\times nK}$.

\subsection{BVAR}
Sea el modelo VAR($p$):
\[
Y_t=C+\Phi_1 Y_{t-1}+\Phi_2 Y_{t-2}+\cdots+\Phi_p Y_{t-p}+E_t,\qquad E_t\overset{i.i.d.}{\sim}\N(0,\Omega),
\]
con $Y_t$ de $n\times1$, $\Phi_j$ de $n\times n$, $E_t$ de $n\times1$, $C$ de $n\times1$ y $t=1,\dots,T$. Definimos
\[
X_t=\begin{bmatrix}1 & Y_{t-1}' & Y_{t-2}' & \cdots & Y_{t-p}'\end{bmatrix}\ \ (1+np\ \text{elementos}=K),\qquad
B=\begin{bmatrix}C & \Phi_1 & \Phi_2 & \cdots & \Phi_p\end{bmatrix}'.
\]
Verificamos que $X_t,B$ reproducen la ecuación original:
\[
X_t B=1\cdot C'+Y_{t-1}'\Phi_1+\cdots+Y_{t-p}'\Phi_p=C+\Phi_1 Y_{t-1}+\cdots+\Phi_p Y_{t-p}.
\]
Apilando $T$ períodos, $Y_t'=X_t B$ para cada $t$:
\[
\underbrace{\begin{bmatrix}Y_1'\\ Y_2'\\ \vdots\\ Y_T'\end{bmatrix}}_{T\times n}
=\underbrace{\begin{bmatrix}X_1\\ X_2\\ \vdots\\ X_T\end{bmatrix}}_{T\times K}
\underbrace{B}_{K\times n}+\underbrace{E}_{T\times n},
\qquad Y=XB+E.
\]
Vectorizando,
\[
\vecop(Y)=(I_n\kron X)\vecop(B)+\vecop(E),\qquad \vecop(E)\sim\N(0,\Omega\kron I_T).
\]

\subsection{Verosimilitud del BVAR}
En el modelo $\by=(I_n\kron X)\beta+\beps$, con $\beps\sim\N(0,\Omega\kron I_T)$, entonces $\by\mid\beta,\Omega\sim\N\big((I_n\kron X)\beta,\ \Omega\kron I_T\big)$, y
\[
\mathcal{L}(\beta,\Omega)=\frac{1}{(2\pi)^{M}\det(\Omega\kron I_T)^{1/2}}
\exp\!\left(-\tfrac12\big(\by-(I_n\kron X)\beta\big)'(\Omega^{-1}\kron I_T)\big(\by-(I_n\kron X)\beta\big)\right).
\]

% ============================================================
\section{Modelo Normal: media desconocida, varianza conocida}
% ============================================================
\textbf{Modelo:} $Y_1,\dots,Y_n\overset{i.i.d.}{\sim}\N(\mu,\sigma^2)$, $i=1,\dots,n$, con $\mu$ desconocida y $\sigma^2$ conocida.\\
\textbf{Problema:} $\mu\mid Y_1,\dots,Y_n=?$

\paragraph{Verosimilitud.}
\[
p(Y_1,\dots,Y_n\mid\mu)=\prod_{i=1}^n\frac{1}{\sqrt{2\pi\sigma^2}}\exp\!\left(-\frac{(Y_i-\mu)^2}{2\sigma^2}\right)
=(2\pi\sigma^2)^{-n/2}\exp\!\left(-\frac{1}{2\sigma^2}\sum_{i=1}^n(Y_i-\mu)^2\right).
\]

\paragraph{Prior.} $\mu\sim\N(\mu_0,\tau_0^2)$.

\paragraph{Posterior.}
\[
p(\mu\mid Y)=\frac{p(Y\mid\mu)\,p(\mu)}{p(Y_1,\dots,Y_n)}
\propto\exp\!\left(-\frac{1}{2\sigma^2}\sum_{i=1}^n(Y_i-\mu)^2\right)\exp\!\left(-\frac{1}{2\tau_0^2}(\mu-\mu_0)^2\right).
\]
Solo mantenemos los términos que dependen de $\mu$. El exponente es
\[
-\tfrac12\left\{\frac{1}{\sigma^2}\Big[\sum_i Y_i^2-2\mu\sum_i Y_i+n\mu^2\Big]+\frac{1}{\tau_0^2}\big[\mu^2-2\mu\mu_0+\mu_0^2\big]\right\}.
\]
Factorizando en $\mu$:
\[
\left(\frac{n}{\sigma^2}+\frac{1}{\tau_0^2}\right)\mu^2-2\left(\frac{\sum_i Y_i}{\sigma^2}+\frac{\mu_0}{\tau_0^2}\right)\mu.
\]
Llamamos $a=\dfrac{n}{\sigma^2}+\dfrac{1}{\tau_0^2}$ y $b=\dfrac{\sum_{i=1}^n Y_i}{\sigma^2}+\dfrac{\mu_0}{\tau_0^2}$. Completando cuadrados,
\[
a\mu^2-2b\mu=a\Big(\mu^2-\tfrac{2b}{a}\mu+\tfrac{b^2}{a^2}-\tfrac{b^2}{a^2}\Big)=a\Big(\mu-\tfrac{b}{a}\Big)^2-\tfrac{b^2}{a},
\]
donde $b^2/a$ no depende de $\mu$. Por lo tanto
\[
\exp\!\left(-\tfrac{a}{2}\Big(\mu-\tfrac{b}{a}\Big)^2\right)=\exp\!\left(-\frac{1}{2\,(1/a)}\Big(\mu-\tfrac{b}{a}\Big)^2\right)
\ \Rightarrow\ \mu_n=\tfrac{b}{a},\ \sigma_n^2=\tfrac{1}{a},
\]
que es el kernel de una densidad normal.

\paragraph{Posterior.}
\[
\mu\mid Y_1,\dots,Y_n\sim\N(\mu_n,\sigma_n^2),
\]
donde
\[
\E[\mu\mid Y]=\frac{b}{a}
=\underbrace{\frac{n/\sigma^2}{n/\sigma^2+1/\tau_0^2}}_{\text{peso de datos}}\bar Y+\underbrace{\frac{1/\tau_0^2}{n/\sigma^2+1/\tau_0^2}}_{\text{peso de prior}}\mu_0,
\qquad
\Var[\mu\mid Y]=a^{-1}=\left(\frac{n}{\sigma^2}+\frac{1}{\tau_0^2}\right)^{-1},
\]
usando $\sum_i Y_i=n\bar Y$. A más datos o más confianza a priori, más precisión a posteriori y varianza más chica. Notar que $\sigma^2/n$ es la varianza de $\bar X$ ($n/\sigma^2$ su precisión) y $\tau_0^2$ es la varianza a priori ($\tau_0^{-2}$ su precisión). \emph{Si el kernel de $p(\mu\mid Y)$ tiene la forma matemática del kernel de una distribución conocida, son la misma.}

\subsection{Distribución predictiva}
Dado que ya observé $Y$, ¿cómo se distribuye una nueva observación $\tilde Y$? Con $Y_i\overset{i.i.d.}{\sim}\N(\mu,\sigma^2)$, $\mu\sim\N(\mu_0,\sigma_0^2)$ y $\mu\mid Y\sim\N(\mu_n,\sigma_n^2)$:
\[
\tilde Y\mid Y\sim\N(\mu_n,\ \sigma^2+\sigma_n^2).
\]
\begin{align*}
\E[\tilde Y\mid Y]&=\E\big[\E[\tilde Y\mid\mu,Y]\big]=\E[\mu\mid Y]=\mu_n,\\
\Var[\tilde Y\mid Y]&=\E\big[\Var[\tilde Y\mid\mu,Y]\big]+\Var\big[\E[\tilde Y\mid\mu,Y]\big]=\sigma^2+\Var[\mu\mid Y]=\sigma^2+\sigma_n^2.
\end{align*}

% ============================================================
\section{Modelo Normal: media conocida, varianza desconocida}
% ============================================================
\textbf{Modelo:} $Y_1,\dots,Y_n\overset{i.i.d.}{\sim}\N(\mu,\sigma^2)$, con $\mu$ conocida y $\sigma^2$ desconocida.\\
\textbf{Problema:} $\sigma^2\mid Y_1,\dots,Y_n$.

\paragraph{Posterior.}
\[
p(\sigma^2\mid Y,\mu)=\frac{p(Y\mid\mu,\sigma^2)\,p(\sigma^2)}{p(Y_1,\dots,Y_n)}
\propto(\sigma^{-2})^{n/2}\exp\!\left(-\sigma^{-2}\sum_{i=1}^n\frac{(Y_i-\mu)^2}{2}\right)p(\sigma^2).
\]

\paragraph{Prior conjugada.} La precisión $\sigma^{-2}\sim\Ga(\alpha,\beta)$, equivalentemente $\sigma^2\sim\IG(\alpha,\beta)$. Con $X\sim\IG(\alpha,\beta)$:
\[
f(x)=\frac{\beta^\alpha}{\Gamma(\alpha)}x^{-\alpha-1}\exp(-\beta/x),\ x>0,
\]
\[
\E[X]=\frac{\beta}{\alpha-1},\qquad \Var[X]=\frac{\beta^2}{(\alpha-1)^2(\alpha-2)},\qquad \text{Moda}[X]=\frac{\beta}{\alpha+1}.
\]

\paragraph{Posterior.}
\begin{align*}
p(\sigma^2\mid Y,\mu)&\propto(\sigma^{-2})^{n/2}\exp\!\left(-\sigma^{-2}\sum_{i=1}^n\frac{(Y_i-\mu)^2}{2}\right)\times(\sigma^2)^{-\alpha-1}\exp(-\beta\sigma^{-2})\\
&=(\sigma^2)^{-0.5n-\alpha-1}\exp\!\left(-\Big(\sum_{i=1}^n\frac{(Y_i-\mu)^2}{2}+\beta\Big)\Big/\sigma^2\right),
\end{align*}
que es el kernel de una Inverse-Gamma. Entonces
\[
\sigma^2\mid Y\sim\IG\!\left(0.5n+\alpha,\ 0.5\sum_{i=1}^n(Y_i-\mu)^2+\beta\right),
\qquad
\E[\sigma^2\mid Y]=\frac{\tilde\beta}{\tilde\alpha-1}=\frac{\tfrac12\sum(Y_i-\mu)^2+\beta}{\tfrac{n}{2}+\alpha-1}.
\]

% ============================================================
\section{Modelo Normal: media y varianza desconocidas}
% ============================================================
\textbf{Modelo:} $Y_1,\dots,Y_n\overset{i.i.d.}{\sim}\N(\mu,\sigma^2)$, con $\mu$ y $\sigma^2$ desconocidas.\\
\textbf{Problema:} $(\mu,\sigma^2)\mid Y_1,\dots,Y_n=?$
\[
p(\mu,\sigma^2\mid Y)\propto p(Y\mid\sigma^2,\mu)\,p(\mu\mid\sigma^2)\,\underbrace{p(\sigma^2)}_{\text{marginal de }\sigma^2}.
\]
Si $\sigma^2$ es conocido, la prior conjugada de $\mu$ es $\N(\mu_0,\tau_0^2)$; en particular asumimos $\tau_0^2=\dfrac{\sigma^2}{\kappa_0}$. Si $\mu$ es conocido, la prior conjugada de $\sigma^2$ es $\IG(\alpha,\beta)$.

\paragraph{Posterior para $\mu$.}
\[
\mu\mid Y_1,\dots,Y_n,\sigma^2\sim\N(\mu_n,\sigma_n^2),
\]
\[
\E[\mu\mid Y,\sigma^2]=\frac{n/\sigma^2}{n/\sigma^2+1/\tau_0^2}\,\bar Y+\frac{1/\tau_0^2}{n/\sigma^2+1/\tau_0^2}\,\mu_0,
\qquad
\Var[\mu\mid Y,\sigma^2]=\left(\frac{n}{\sigma^2}+\frac{\kappa_0}{\sigma^2}\right)^{-1}=\frac{\sigma^2}{\kappa_0+n}.
\]

\paragraph{Posterior para $\sigma^2$ (marginalizando $\mu$).}
\[
p(\sigma^2\mid Y)\propto p(\sigma^2)\int p(Y_1,\dots,Y_n\mid\mu,\sigma^2)\,p(\mu\mid\sigma^2)\,\dpar\mu.
\]
Con verosimilitud $\propto(\sigma^2)^{-n/2}\exp\!\big(-\tfrac{\sum(Y_i-\mu)^2}{2\sigma^2}\big)$, prior de $\mu\mid\sigma^2$ con $\tau_0^2=\sigma^2/\kappa_0$ $\propto(\sigma^2)^{-1/2}\exp\!\big(-\tfrac{\kappa_0(\mu-\mu_0)^2}{2\sigma^2}\big)$, y prior de $\sigma^2$ $\propto(\sigma^2)^{-\alpha-1}\exp(-\beta/\sigma^2)$. Multiplicando verosimilitud y prior de $\mu$ y juntando exponentes:
\[
(\sigma^2)^{-(n+1)/2}\int\exp\!\left(-\frac{1}{2\sigma^2}\Big[\sum(Y_i-\mu)^2+\kappa_0(\mu-\mu_0)^2\Big]\right)\dpar\mu.
\]
Expandiendo,
\[
\sum(Y_i-\mu)^2+\kappa_0(\mu-\mu_0)^2=(n+\kappa_0)\mu^2-2(n\bar Y+\kappa_0\mu_0)\mu+\underbrace{\big(\textstyle\sum Y_i^2+\kappa_0\mu_0^2\big)}_{\text{no depende de }\mu}.
\]
Con $A=n+\kappa_0$ y $B=n\bar Y+\kappa_0\mu_0$, completando cuadrados y notando que $B/A=\mu_n$,
\[
A\mu^2-2B\mu=A\Big(\mu-\tfrac{B}{A}\Big)^2-\tfrac{B^2}{A}=(n+\kappa_0)(\mu-\mu_n)^2-\frac{(n\bar Y+\kappa_0\mu_0)^2}{n+\kappa_0}.
\]
La integral en $\mu$ es
\[
\int\exp\!\left(-\frac{(n+\kappa_0)(\mu-\mu_n)^2}{2\sigma^2}\right)\dpar\mu=\sqrt{2\pi\,\frac{\sigma^2}{n+\kappa_0}}\ \propto\ (\sigma^2)^{1/2},
\]
usando que toda densidad normal integra $1$. Reemplazando,
\begin{multline*}
(\sigma^2)^{-n/2}\cdot(\sigma^2)^{-1/2}\cdot(\sigma^2)^{1/2}\cdot(\sigma^2)^{-\alpha-1}\\
=(\sigma^2)^{-n/2-\alpha-1}\exp\!\left(-\frac{1}{\sigma^2}\underbrace{\left[\frac{\sum_i Y_i^2+\kappa_0\mu_0^2}{2}-\frac{(n\bar Y+\kappa_0\mu_0)^2}{2(n+\kappa_0)}+\beta\right]}_{\tilde\beta}\right).
\end{multline*}
Comparando exponentes: $-n/2-\alpha-1=-\tilde\alpha-1\Rightarrow\tilde\alpha=\alpha+n/2$; y $\tilde\beta/\sigma^2$ ya identificado. Por lo tanto
\[
\sigma^2\mid Y\sim\IG(\tilde\alpha,\tilde\beta),\qquad
\mu\mid Y_1,\dots,Y_n,\sigma^2\sim\N(\mu_n,\sigma_n^2).
\]
Se desea entender $\E[\mu\mid Y]$, $\Var[\mu\mid Y]$ y $P(\mu_1>\mu_2\mid Y)$. No se tiene la marginal de $\mu$ explícita.

\subsection{Aproximación de Monte Carlo}
\begin{enumerate}[leftmargin=2em]
  \item Simular $\sigma^{2(i)}\sim\IG$.
  \item Simular $\mu^{(i)}\sim\N$.
\end{enumerate}
Entonces $(\sigma^{2(1)},\mu^{(1)}),\dots,(\sigma^{2(S)},\mu^{(S)})$ es una muestra de la distribución conjunta a posteriori.

% ============================================================
\section{Algoritmo de Gibbs Sampling}
% ============================================================

\subsection{Recordar: Monte Carlo vs.\ Gibbs}
El algoritmo de Gibbs trata de aproximar una distribución $f(X,Y)$.

\paragraph{¿Cómo? Monte Carlo.} Se factoriza
\[
f(X,Y)=\underbrace{f(Y\mid X)}_{Y\mid X}\underbrace{f(X)}_{\text{modelo marginal de }X}.
\]
Con estos dos submodelos (\emph{independientes}) se generan pares independientes:
\[
X_1\to Y_1\Rightarrow(X_1,Y_1),\quad X_2\to Y_2\Rightarrow(X_2,Y_2),\quad X_3\to Y_3\Rightarrow(X_3,Y_3),\ \dots
\]
y se realizan las aproximaciones de Monte Carlo. Por la \textbf{Ley de los Grandes Números} (los promedios de la muestra convergen a los promedios del modelo):
\[
\E(XY)\xleftarrow{\ p\ }\frac{1}{n}\sum_{i=1}^n X_i Y_i,\qquad
\E(X)\xleftarrow{\ p\ }\frac{1}{n}\sum_i X_i,\qquad
\Cov(X,Y)\xleftarrow{\ p\ }\hat\rho.
\]
Requisito: errores i.i.d.

\paragraph{Gibbs Sampling.} Permite las aproximaciones \emph{con dependencia} en la muestra. Necesito las condicionales $f(X\mid Y)$ y $f(Y\mid X)$, con las que genero una secuencia dependiente. Comienzo con ``$X_1$'' (viene del modelo $f(X)$ aleatoriamente), luego ``$Y_1$'' $\to$ ``$X_2$'' $\to$ ``$Y_2$'' $\to\cdots$: me turno entre las distribuciones condicionales. Se generan pares dependientes
\[
(X_1,Y_1),\ (X_2,Y_2),\ (X_3,Y_3),\ \dots
\]
Toda la información del pasado se guarda dentro de la muestra ``actual'': \textbf{es una cadena de Markov}.
\begin{itemize}[leftmargin=2em]
  \item No se puede aplicar la Ley de los Grandes Números (las muestras no son independientes).
  \item En general, \emph{aproxima} una muestra independiente.
  \item En Inferencia Bayesiana, el objetivo es aproximar $p(\theta\mid\text{Data})$.
\end{itemize}

\subsection{Ejemplo}
Mismo modelo $Y_1,\dots,Y_n\sim\N(\mu,\sigma^2)$. La prior ahora se trabaja con $\sigma$ (desviación estándar):
\[
\mu\mid\sigma^2\sim\N\!\left(\mu_0,\tfrac{\sigma}{\sqrt{\kappa_0}}\right),\qquad
1/\sigma\sim\Ga\!\left(\tfrac{\nu_0}{2},\tfrac{\nu_0\sigma_0^2}{2}\right).
\]

\paragraph{Algoritmo.}
\begin{itemize}[leftmargin=2em]
  \item Genera $\sigma^{2(i)}$ usando $p(\sigma^2\mid Y_1,\dots,Y_n)$.
  \item Genera $\mu^{(i)}$ usando $p(\mu\mid\sigma^{2(i)},Y_1,\dots,Y_n)$.
\end{itemize}
La conjunta a posteriori $(\mu^{(1)},\sigma^{2(1)}),\dots,(\mu^{(S)},\sigma^{2(S)})$ es una muestra aleatoria de $p(\mu,\sigma^2\mid Y)$; la marginal $(\mu^{(1)},\dots,\mu^{(S)})$ es una muestra de $p(\mu\mid Y)$.

\paragraph{Problema.} Ya no es posible establecer explícitamente la posterior. Sabemos que
\[
p(\mu,\sigma^2\mid Y_1,\dots,Y_n)=\frac{p(Y_1,\dots,Y_n\mid\mu,\sigma^2)\,p(\mu,\sigma^2)}{p(Y_1,\dots,Y_n)}.
\]
Podríamos aproximar de manera discreta $p(Y_1,\dots,Y_n)$. Sea $\{\mu_1,\dots,\mu_s\}\in R_\mu$ y $\{\sigma_1^2,\dots,\sigma_s^2\}\in R_{\sigma^2}$; entonces
\[
p(\mu_i,\sigma_i^2\mid Y_1,\dots,Y_n)=\frac{p(Y\mid\sigma_j^2,\mu_i)\,p(\mu_i,\sigma_i^2)}{\sum_{i=1}^s\sum_{j=1}^T p(\mu_i,\sigma_j^2,Y_1,\dots,Y_n)}
\]
aproxima la posterior, pero cae en la \textbf{maldición de la dimensionalidad}.

\subsection{Solución vía condicionales completas}
\paragraph{$p(\sigma^2\mid\mu,Y)$.}
\begin{align*}
p(\sigma^2\mid\mu,Y)&\propto p(Y\mid\sigma^2,\mu)\,p(\mu\mid\sigma^2)\,p(\sigma^2)\\
&\propto\sigma^{-n/2}\exp\!\left(-\sigma^{-2}\sum_{i=1}^n(Y_i-\mu)^2/2\right)\times(\sigma^2)^{-\alpha-1}\exp(-\beta\sigma^{-2})\\
&=(\sigma^{-2})^{n/2+\alpha+1}\exp\!\left(-\sigma^{-2}\Big(\sum_{i=1}^n(Y_i-\mu)^2/2+\beta\Big)\right),
\end{align*}
(el término $p(\mu\mid\sigma^2)$ se ignora por no depender de $\sigma^2$ una vez fijado $\mu$). Entonces
\[
\sigma^{-2}\mid\mu,Y\sim\Ga\!\left(0.5n+\alpha,\ 0.5\sum_{i=1}^n(Y_i-\mu)^2+\beta\right).
\]
Análogamente se sabe $p(\mu\mid\sigma^2,Y)$.

\paragraph{Esquema secuencial (Markov).} Supuesto: se tiene $\sigma^{2(0)}$ generado de $p(\sigma^2\mid Y)$.
\begin{enumerate}[leftmargin=2em]
  \item Generar $\mu^{(1)}$ usando $p(\mu\mid\sigma^{2(0)},\bar Y)$.
  \item Generar $\sigma^{2(1)}$ usando $p(\sigma^2\mid\mu^{(1)},\bar Y)$.
\end{enumerate}
Repitiendo se genera la secuencia dependiente $\sigma^{2(0)},\mu^{(0)},\sigma^{2(1)},\mu^{(1)},\dots$ Esta secuencia \emph{no} es una muestra de Monte Carlo.

\subsection{Gibbs Sampling general}
Sea $\phi=(\phi_1,\dots,\phi_p)$. Insumos: prior $\phi\sim p(\phi)$ y punto de inicio $\phi^{(0)}=(\phi_1^{(0)},\dots,\phi_p^{(0)})$. Se genera $\phi^{(i)}$ vía
\begin{align*}
\phi_1^{(i)}&\sim p\big(\phi_1\mid\phi_2^{(i-1)},\dots,\phi_p^{(i-1)},Y\big),\\
\phi_2^{(i)}&\sim p\big(\phi_2\mid\phi_1^{(i)},\phi_3^{(i-1)},\dots,\phi_p^{(i-1)},Y\big),\\
\phi_3^{(i)}&\sim p\big(\phi_3\mid\phi_1^{(i)},\phi_2^{(i)},\phi_4^{(i-1)},\dots,\phi_p^{(i-1)},Y\big),\\
&\ \ \vdots\\
\phi_p^{(i)}&\sim p\big(\phi_p\mid\phi_1^{(i)},\phi_2^{(i)},\dots,\phi_{p-1}^{(i)},Y\big).
\end{align*}
En general se deben conocer todas las distribuciones condicionales. Se genera una secuencia de vectores dependientes $\phi^{(1)},\dots,\phi^{(S)}$.

\paragraph{Propiedad de Markov.} $\phi^{(i)}$ depende de $\phi^{(0)},\dots,\phi^{(i-1)}$ solo a través de $\phi^{(i-1)}$: el futuro depende solamente del pasado. Esto es una \textbf{cadena de Markov}. Bajo ciertas condiciones,
\[
P\big(\phi^{(S)}\in A\mid Y\big)\xrightarrow{S\to\infty}\int_A p(\phi\mid Y)\,\dpar\phi,
\]
es decir, la distribución de $\phi^{(S)}$ converge a la posterior. Además el promedio de la cadena converge al valor esperado a posteriori:
\[
\frac{1}{S}\sum_{s=1}^S g(\phi^{(s)})\xrightarrow{S\to\infty}\E[g(\phi)\mid Y].
\]

\subsection{Los dos problemas del MCMC}
\begin{itemize}[leftmargin=2em]
  \item \textbf{Estacionaridad:} la cadena necesita olvidar su punto de inicio y llegar a una región representativa de la posterior. Se soluciona empezando en una región de alta probabilidad, verificando que dos submuestras de la cadena tengan la misma distribución, o usando muchas iteraciones y descartando las primeras (\emph{burn-in}).
  \item \textbf{Mixing (velocidad de mezcla):} incluso cuando la cadena ya convergió, puede moverse muy lento entre distintas regiones de alta probabilidad. Esto genera alta correlación y reduce la eficiencia de la estimación.
\end{itemize}

\subsection{Varianza del estimador Monte Carlo}
Con muestra independiente $\phi^{(1)},\dots,\phi^{(S)}$ y $\bar\phi=\tfrac{1}{S}\sum_{s=1}^S\phi^{(s)}$:
\[
\Var_{\text{MC}}[\bar\phi]=\E\big[(\bar\phi-\E[\bar\phi])^2\big]=\frac{\Var(\phi)}{S},
\qquad
\text{IC}_{95\%}\text{ para }\E[\bar\phi]:\ \ \bar\phi\pm 1.96\sqrt{\tfrac{\Var(\phi)}{S}}.
\]

\subsection{Varianza MCMC}
\begin{align*}
\Var_{\text{MCMC}}[\bar\phi]&=\E\big[(\bar\phi-\phi_0)^2\big]=\E\!\left[\Big(\tfrac{1}{S}\sum_{s=1}^S(\phi^{(s)}-\phi_0)\Big)^2\right]\\
&=\frac{1}{S^2}\E\!\left[\sum_{s=1}^S(\phi^{(s)}-\phi_0)^2+\sum_{s\ne t}(\phi^{(s)}-\bar\phi)(\phi^{(t)}-\bar\phi)\right]\\
&=\Var_{\text{MC}}[\bar\phi]+\frac{1}{S^2}\sum_{s\ne t}\Cov(\phi^{(s)},\phi^{(t)}),
\end{align*}
donde el segundo término es positivo cuando las muestras están correlacionadas. La autocorrelación empírica de la cadena es
\[
\operatorname{acf}_\tau(\phi)=\frac{\tfrac{1}{S-\tau}\sum_{s=1}^{S-\tau}(\phi^{(s)}-\bar\phi)(\phi^{(s+\tau)}-\bar\phi)}{\tfrac{1}{S-1}\sum_{s=1}^S(\phi^{(s)}-\bar\phi)^2}.
\]

\paragraph{Tamaño de muestra efectiva.}
\[
\Var_{\text{MCMC}}[\bar\phi]=\frac{\Var(\phi)}{S_{\text{eff}}},
\]
donde $S_{\text{eff}}$ es el número de muestras independientes que darían la misma precisión que la muestra MCMC obtenida.

% ============================================================
\section{Ejemplo: número de accidentes (punto de cambio de Poisson)}
% ============================================================
\textbf{Modelo:}
\[
X_i\sim\text{Poisson}(\mu),\ i=1,\dots,k;\qquad X_i\sim\text{Poisson}(\lambda),\ i=k+1,\dots,m.
\]
Parámetros $(\mu,\lambda,k)$. \textbf{Priors:} $\mu\sim\Ga(\alpha_1,\beta_1)$, $\lambda\sim\Ga(\alpha_2,\beta_2)$, $k\sim\text{Uniforme}\{1,2,\dots,m\}$.

\paragraph{Verosimilitud.}
\[
p(\underline X\mid\mu,\lambda,k)=\prod_{i=1}^k\frac{\mu^{X_i}e^{-\mu}}{X_i!}\prod_{i=k+1}^m\frac{\lambda^{X_i}e^{-\lambda}}{X_i!}.
\]

\paragraph{Posterior.}
\begin{align*}
p(\mu,\lambda,k\mid\underline X)&\propto p(\underline X\mid\mu,\lambda,k)\,p(\mu)\,p(\lambda)\,\underbrace{p(k)}_{1/m}\\
&\propto\mu^{\sum_{i=1}^k X_i}e^{-\mu k}\lambda^{\sum_{i=k+1}^m X_i}e^{-\lambda(m-k)}\times\mu^{\alpha_1-1}e^{-\beta_1\mu}\lambda^{\alpha_2-1}e^{-\beta_2\lambda}\tfrac{1}{m}\\
&=\mu^{\alpha_1+\sum_{i=1}^k X_i-1}e^{-\mu(k+\beta_1)}\ \lambda^{\alpha_2+\sum_{i=k+1}^m X_i-1}e^{-\lambda(m-k+\beta_2)}.
\end{align*}

\paragraph{Condicionales de $\mu$ y $\lambda$.}
\[
\mu\mid\underline X,\lambda,k\sim\Ga\!\left(\alpha_1+\sum_{i=1}^k X_i,\ k+\beta_1\right),\qquad
\lambda\mid\underline X,\mu,k\sim\Ga\!\left(\alpha_2+\sum_{i=k+1}^m X_i,\ (m-k)+\beta_2\right).
\]

\paragraph{Condicional de $k$.} Fijando $\mu$ y $\lambda$, los términos que dependen de $k$ son
\[
p(k\mid\underline X,\mu,\lambda)\propto\mu^{\sum_{i=1}^k X_i}\,\lambda^{\sum_{i=k+1}^m X_i}\,e^{-k(\mu-\lambda)}=\mathcal{L}(\underline X,\lambda,\mu,k).
\]
Como $k$ es discreto (el punto de cambio es un mes en particular) y se distribuye uniformemente,
\[
p(k\mid\mu,\lambda,\underline X)=\frac{\mathcal{L}(\underline X,\lambda,\mu,k)}{\sum_{i=1}^m\mathcal{L}(\underline X,\lambda,\mu,i)},\ i=1,\dots,m,
\]
(numéricamente se aplican logaritmos y se resta el máximo).

% ============================================================
\section{Regresión lineal (parametrización con precisión)}
% ============================================================
\textbf{Modelo:} $Y_i=\beta_0+\beta_1 X_i+\varepsilon_i$, $\varepsilon_i\sim\N(0,1/\tau)$.\\
\textbf{Priors:} $\beta_0\sim\N(0,1/\tau_0)$, $\beta_1\sim\N(0,1/\tau_1)$, $\tau\sim\Ga(\alpha,\beta)$.

\paragraph{Verosimilitud.}
\[
p(\underline Y\mid\beta_0,\beta_1,\tau)=\left(\frac{1}{\sqrt{2\pi}}\right)^n\tau^{n/2}\exp\!\left(-\frac{\tau}{2}\sum_{i=1}^n(Y_i-\beta_0-\beta_1 X_i)^2\right).
\]

\paragraph{Posterior.}
\begin{align*}
p(\beta_0,\beta_1,\tau\mid\underline Y,\underline X)
&=p(\underline Y\mid\beta_0,\beta_1,\tau)\,p(\beta_0)\,p(\beta_1)\,p(\tau)\\
&=\tau^{n/2}\exp\!\left(-\tfrac{\tau}{2}\sum_{i=1}^n(Y_i-\beta_0-\beta_1 X_i)^2\right)e^{-0.5\tau_0\beta_0^2}e^{-0.5\tau_1\beta_1^2}\tau^{\alpha-1}e^{-\beta\tau}.
\end{align*}

\paragraph{Condicional de $\tau$.}
\[
p(\tau\mid\beta_0,\beta_1)\propto\tau^{0.5n+\alpha-1}\exp\!\left(-\tau\Big[0.5\sum_{i=1}^n(Y_i-\beta_0-\beta_1 X_i)^2+\beta\Big]\right),
\]
entonces
\[
\tau\mid\beta_0,\beta_1\sim\Ga\!\left(0.5n+\alpha,\ 0.5\sum_{i=1}^n(Y_i-\beta_0-\beta_1 X_i)^2+\beta\right).
\]

\paragraph{Condicional de $\beta_0$.}
\[
p(\beta_0\mid\tau,\beta_1)\propto\exp\!\left(-0.5\tau\sum_{i=1}^n\big[(Y_i-\beta_1 X_i)-\beta_0\big]^2-0.5\tau_0\beta_0^2\right).
\]
Con $r_i=Y_i-\beta_1 X_i$ y $\bar r=\bar Y-\beta_1\bar X$:
\[
\sum(r_i-\beta_0)^2=\sum r_i^2-2\beta_0\sum r_i+n\beta_0^2,
\]
(el primer término se omite). El exponente queda $-\tfrac{\tau}{2}[-2\beta_0 n\bar r+n\beta_0^2]-0.5\tau_0\beta_0^2=-0.5[(\tau n+\tau_0)\beta_0^2-2\tau n\bar r\beta_0]$. Con $A=\tau n+\tau_0$ y $B=\tau n\bar r$, completando cuadrados,
\[
A\beta_0^2-2B\beta_0=A\Big(\beta_0-\tfrac{B}{A}\Big)^2-\tfrac{B^2}{A}\ \Rightarrow\ -\tfrac{A}{2}\Big(\beta_0-\tfrac{B}{A}\Big)^2=-\frac{1}{2\,(1/A)}\Big(\beta_0-\tfrac{B}{A}\Big)^2.
\]
Identificando la normal,
\[
\beta_0\mid\beta_1,\tau\sim\N\!\left(\frac{\tau n(\bar Y-\beta_1\bar X)}{\tau n+\tau_0},\ \frac{1}{\tau n+\tau_0}\right).
\]

\paragraph{Condicional de $\beta_1$.}
\[
p(\beta_1\mid\tau,\beta_0)\propto\exp\!\left(-\tfrac{\tau}{2}\sum_{i=1}^n(Y_i-\beta_0-\beta_1 X_i)^2\right)\exp\!\left(-\tfrac{\tau_1}{2}\beta_1^2\right).
\]
Con $u_i=Y_i-\beta_0$: $\sum(u_i-\beta_1 X_i)^2=\sum u_i^2-2\beta_1\sum u_i X_i+\beta_1^2\sum X_i^2$. Usando $\sum X_i^2=n\bar X^2$ y $\sum u_i X_i=n\overline{XY}-n\bar X\beta_0$:
\[
-\tfrac12\big[(\tau n\bar X^2+\tau_1)\beta_1^2-2\tau(n\overline{XY}-n\bar X\beta_0)\beta_1\big].
\]
Con $A=\tau n\bar X^2+\tau_1$ y $B=\tau(n\overline{XY}-n\bar X\beta_0)$,
\[
\beta_1\mid\beta_0,\tau\sim\N\!\left(\frac{\tau(n\overline{XY}-n\bar X\beta_0)}{\tau n\bar X^2+\tau_1},\ \frac{1}{\tau n\bar X^2+\tau_1}\right).
\]

% ============================================================
\section{Modelo Beta-Binomial}
% ============================================================
\textbf{Modelo:} $X\mid\theta\sim\text{Binomial}(n,\theta)$, $\theta\sim\text{Beta}(\alpha,\beta)$. Modelo conjunto
\[
f(X,\theta)=\binom{n}{X}\theta^{X}(1-\theta)^{n-X}\cdot\frac{\theta^{\alpha-1}(1-\theta)^{\beta-1}}{B(\alpha,\beta)}.
\]

% ============================================================
\section{Mezcla de normales con variable latente}
% ============================================================
\textbf{Modelo para $X$:}
\[
f(X\mid\pi,\mu_0,\mu_1)=(1-\pi)\,\N(X\mid\mu_0,\tau^{-1})+\pi\,\N(X\mid\mu_1,\tau^{-1}),
\]
donde $\tau$ es la precisión (conocida). Hay 3 parámetros desconocidos $(\pi,\mu_0,\mu_1)$ con priors
\[
\pi\sim\text{Beta}(\alpha,\beta)\ (0\le\pi\le1),\qquad \mu_0\sim\N(\eta_0,\tau_0),\qquad \mu_1\sim\N(\eta_1,\tau_1).
\]
¿Cuál es la posterior $p(\pi,\mu_0,\mu_1\mid X)$? La verosimilitud tiene una suma dentro del producto, lo que dificulta calcular la posterior. \textbf{Acudimos a variables auxiliares (latentes).}

\paragraph{Variable latente.}
\[
Z_i=\begin{cases}0 & \text{si }X_i\text{ viene del grupo 0 (media }\mu_0)\\ 1 & \text{si }X_i\text{ viene del grupo 1 (media }\mu_1)\end{cases}
\qquad Z_i\mid\pi\sim\text{Bernoulli}(\pi).
\]
Entonces
\begin{align*}
p(X_i\mid\mu,\pi)&=p(X_i\mid Z_i=0,\mu,\pi)P(Z_i=0\mid\mu,\pi)+p(X_i\mid Z_i=1,\mu,\pi)P(Z_i=1\mid\mu,\pi)\\
&=(1-\pi)\N(X_i\mid\mu_0,\tau^{-1})+\pi\N(X_i\mid\mu_1,\tau^{-1})=f(X\mid\mu,\pi),
\end{align*}
con $\mu=(\mu_0,\mu_1)$, y $X_i\mid Z_i=0\sim\N(\mu_0,\tau^{-1})$, $X_i\mid Z_i=1\sim\N(\mu_1,\tau^{-1})$. Ahora tenemos un producto de verosimilitudes:
\[
p(X\mid Z,\mu,\pi)=\prod_{Z_i=0}\N(X_i\mid\mu_0,\tau^{-1})\prod_{Z_i=1}\N(X_i\mid\mu_1,\tau^{-1}).
\]

\paragraph{Posterior.}
\[
p(\mu,\pi,\underline Z\mid X)\propto\prod_{i=1}^n p(X_i\mid Z_i,\mu,\pi)\,p(Z_i\mid\pi,\mu)\,p(\pi,\mu),
\]
usando la regla de la cadena $p(A,B,C)=p(A\mid B,C)p(B\mid C)p(C)$ con $A=X$, $B=Z$, $C=(\mu,\pi)$, y $p(\mu,\pi)=p(\mu_1)p(\mu_0)p(\pi)$.

\paragraph{Condicional de $\mu_0$ (análogo $\mu_1$).}
\[
p(\mu_0\mid\pi,\underline Z,X)\propto\prod_{Z_i=0}p(X_i\mid\mu_0)\,p(\mu_0)\propto\exp\!\left(-\tfrac{\tau}{2}\sum_{Z_i=0}(X_i-\mu_0)^2\right)\exp\!\left(-\tfrac{\tau_0}{2}(\mu_0-\eta_0)^2\right).
\]
Con $(X_i-\mu_0)^2=X_i^2-2X_i\mu_0+\mu_0^2$, sumando sobre $Z_i=0$ ($n_0$ observaciones, media $\bar X_0$):
\[
\sum_{Z_i=0}(X_i-\mu_0)^2\propto-2\mu_0 n_0\bar X_0+n_0\mu_0^2.
\]
De la prior $(\mu_0-\eta_0)^2\propto\mu_0^2-2\mu_0\eta_0$. Metiendo todo dentro,
\[
p(\mu_0\mid\cdots)\propto\exp\!\left(-\tfrac12\big[(n_0\tau+\tau_0)\mu_0^2-2(n_0\tau\bar X_0+\tau_0\eta_0)\mu_0\big]\right).
\]
Con $A=n_0\tau+\tau_0$ y $B=n_0\tau\bar X_0+\tau_0\eta_0$, completando cuadrados,
\[
\mu_0\mid\cdots\sim\N\!\left(\frac{n_0\tau}{n_0\tau+\tau_0}\,\bar X_0+\frac{\tau_0}{n_0\tau+\tau_0}\,\eta_0,\ (n_0\tau+\tau_0)^{-1}\right).
\]

\paragraph{Condicional de $\mu_1$ (por el capítulo 2).}
\[
\mu_1\mid\pi,\underline Z,\underline X\sim\N\!\left(\frac{n_1\tau}{n_1\tau+\tau_1}\bar X_1+\frac{\tau_1}{n_1\tau+\tau_1}\eta_1,\ (n_1\tau+\tau_1)^{-1}\right),
\quad n_1=\sum_{i=1}^n Z_i,\ \bar X_1=\frac{\sum_i X_i Z_i}{\sum_i Z_i}.
\]

\paragraph{Condicional de $\pi$.}
\[
p(\pi\mid\underline Z,\mu,\underline X)\propto\pi^{n_1+\alpha-1}(1-\pi)^{n_0+\beta-1}\ \Rightarrow\ \pi\mid\underline Z,\mu,\underline X\sim\text{Beta}(n_1+\alpha,\ n_0+\beta).
\]

\paragraph{Condicional de $Z_i$.}
\[
P(Z_i=1\mid\mu,\pi)=\frac{p(x_i\mid z_i=1,\mu,\pi)P(Z_i=1\mid\mu,\pi)}{Pa},
\]
\begin{align*}
Pa&=p(x_i\mid z_i=0,\mu,\pi)P(Z_i=0\mid\mu,\pi)+p(x_i\mid z_i=1,\mu,\pi)P(Z_i=1\mid\mu,\pi)\\
&=(1-\pi)\N(\mu_0,\sqrt{1/\tau})+\pi\N(\mu_1,\sqrt{1/\tau}).
\end{align*}

% ============================================================
\section{Distribución Normal multivariada}
% ============================================================
$X\sim\N_K(\mu,\Sigma)$ tiene densidad
\[
f_X(x)=\frac{1}{|\Sigma|^{1/2}(2\pi)^{K/2}}\exp\!\left(-\tfrac12(x-\mu)^{\top}\Sigma^{-1}(x-\mu)\right),
\]
con $\E(X)=\mu=(\mu_1,\dots,\mu_K)'$ y $\Sigma$ la matriz de covarianzas ($\Var(X_i)=\sigma_i^2$, $\Cov(X_i,X_j)=\sigma_{ij}$). Marginal: $X_i\sim\N(\mu_i,\sigma_i^2)$.
\begin{align*}
f_X(x)&\propto\exp\!\left(-\tfrac12(x-\mu)^{\top}\Sigma^{-1}(x-\mu)\right)
\propto\exp\!\left(-\tfrac12\big[x^{\top}\Sigma^{-1}x-2x^{\top}\Sigma^{-1}\mu+\mu^{\top}\Sigma^{-1}\mu\big]\right)\\
&\propto\exp\!\left(-\tfrac12\big[x^{\top}\underbrace{\Sigma^{-1}}_{A}x-2x^{\top}\underbrace{\Sigma^{-1}\mu}_{a}\big]\right).
\end{align*}

\begin{prop}
Si la densidad de un vector $X$ es $f(x)\propto\exp\!\big(-\tfrac12[x^{\top}Ax-2x^{\top}a]\big)$, con $A$ matriz $K\times K$ positiva definida y $a$ vector $K\times1$, entonces
\[
X\sim\N(A^{-1}a,\ A^{-1}).
\]
\end{prop}
Esto se logra expandiendo el producto de la normal multivariada:
\begin{align*}
\tfrac12(y-\mu)^{\top}\Sigma^{-1}(y-\mu)&=-\tfrac12\big(y^{\top}\Sigma^{-1}y-y^{\top}\Sigma^{-1}\mu-\mu^{\top}\Sigma^{-1}y+\mu^{\top}\Sigma^{-1}\mu\big),
\end{align*}
y como $\Sigma^{-1}$ es simétrica, $y^{\top}\Sigma^{-1}\mu=\mu^{\top}\Sigma^{-1}y$, de modo que
$=-\tfrac12(y^{\top}\Sigma^{-1}y-2y^{\top}\Sigma^{-1}\mu)+\text{cte}$. Comparando: $A=\Sigma^{-1}\Rightarrow\Sigma=A^{-1}$ y $a=\Sigma^{-1}\mu\Rightarrow A^{-1}a=\mu$.

% ============================================================
\section{Modelo de regresión lineal (revisado)}
% ============================================================
\subsection{Distribución Gamma e Inverse-Gamma}
$Y\sim\Ga(a,b)$: $f(Y)=\dfrac{b^a}{\Gamma(a)}Y^{a-1}e^{-bY}$, $Y>0$, $a,b>0$, con $\E[Y]=a/b$, $\Var(Y)=a/b^2$.\\
$Y\sim\IG(a,b)$: $f_Y(Y)=\dfrac{b^a}{\Gamma(a)}\left(\dfrac{1}{Y}\right)^{a+1}e^{-b/Y}$, $Y>0$. Propiedad: $Y\sim\Ga(a,b)\iff 1/Y\sim\IG(a,b)$. Es la prior conjugada para la varianza en modelos normales.

\subsection{Planteamiento}
$y_i=X_i^{\top}\beta+\varepsilon_i$, $\varepsilon_i\overset{ind}{\sim}\N(0,\sigma^2)$, $i=1,\dots,n$, es decir $\by=\bX\beta+\beps$ con $\beps\sim\N(0,\sigma^2\bI)$. Por lo tanto $\by\mid\beta,\sigma^2\sim\N(\bX\beta,\sigma^2\bI)$ y la máxima verosimilitud es
\[
p(\by\mid\beta,\sigma^2)=\left(\frac{1}{2\pi}\right)^{n/2}\left(\frac{1}{\sigma^2}\right)^{n/2}\exp\!\left(-\frac{1}{2\sigma^2}(\by-\bX\beta)^{\top}(\by-\bX\beta)\right).
\]

\subsection{Distribuciones a priori}
\begin{itemize}[leftmargin=2em]
  \item \textbf{No informativa (impropia):} $p(\beta,\sigma^2)\propto\dfrac{1}{\sigma^2}$ (no integra a $1$, pero es válida para la posterior).
  \item \textbf{Informativa:} $p(\beta,\sigma^2)=p(\beta)p(\sigma^2)$ con $\beta\sim\N(\beta_0,\Sigma_0)$ y $\sigma^2\sim\IG(a_0,b_0)$.
\end{itemize}

\subsection{Posterior con prior no informativa}
\[
p(\beta,\sigma^2\mid\by)\propto p(\by\mid\beta,\sigma^2)\,p(\beta,\sigma^2)
\propto\left(\frac{1}{\sigma^2}\right)^{\frac{n}{2}+1}\exp\!\left(-\frac{1}{2\sigma^2}(\by-\bX\beta)^{\top}(\by-\bX\beta)\right).
\]

\subsection{Condicional completa de $\beta$}
\begin{align*}
p(\beta\mid\sigma^2,\by)&\propto\exp\!\left(-\frac{1}{2\sigma^2}\big[\underbrace{\by^{\top}\by}_{\text{sale}}-2\beta^{\top}\bX^{\top}\by+\beta^{\top}\bX^{\top}\bX\beta\big]\right)\\
&\propto\exp\!\left(-\tfrac12\Big[\beta^{\top}\underbrace{\tfrac{\bX^{\top}\bX}{\sigma^2}}_{A}\beta-2\beta^{\top}\underbrace{\tfrac{\bX^{\top}\by}{\sigma^2}}_{a}\Big]\right).
\end{align*}
Comparando con el kernel de una normal multivariada $\propto\exp\!\big(-\tfrac12(\beta^{\top}\Sigma^{-1}\beta-2\beta^{\top}\Sigma^{-1}\mu)\big)$:
\[
\Sigma^{-1}=\frac{\bX^{\top}\bX}{\sigma^2}\Rightarrow\Sigma=\sigma^2(\bX^{\top}\bX)^{-1},\qquad
\Sigma^{-1}\mu=\frac{\bX^{\top}\by}{\sigma^2}\Rightarrow\mu=(\bX^{\top}\bX)^{-1}\bX^{\top}\by=\hat\beta_{\text{OLS}},
\]
de donde
\[
\beta\mid\sigma^2,\by\sim\N\!\left((\bX^{\top}\bX)^{-1}\bX^{\top}\by,\ \sigma^2(\bX^{\top}\bX)^{-1}\right).
\]

\subsection{Condicional completa de $\sigma^2$}
Fijando $\beta$ y con $\text{SSE}=(\by-\bX\beta)^{\top}(\by-\bX\beta)$:
\[
p(\sigma^2\mid\beta,\by)\propto\left(\frac{1}{\sigma^2}\right)^{n/2+1}\exp\!\left(-\frac{\text{SSE}}{2\sigma^2}\right).
\]
Comparando con el kernel de $\IG(a,b)$, $\left(\tfrac{1}{\sigma^2}\right)^{a+1}\exp(-b/\sigma^2)$: $a+1=\tfrac{n}{2}+1\Rightarrow a=\tfrac{n}{2}$ y $b=\tfrac{\text{SSE}}{2}$. Por lo tanto
\[
\sigma^2\mid\beta,\by\sim\IG\!\left(\frac{n}{2},\ \frac{\text{SSE}}{2}\right),\qquad
\text{SSE}=(\by-\bX\beta)^{\top}(\by-\bX\beta).
\]

\subsection{Gibbs Sampling}
\begin{enumerate}[leftmargin=2em]
  \item Simular $\beta^{t+1}\sim\N\big((\bX^{\top}\bX)^{-1}\bX^{\top}\by,\ \sigma_t^2(\bX^{\top}\bX)^{-1}\big)$.
  \item Simular $\sigma^2_{t+1}\sim\IG\big(\tfrac{n}{2},\ \tfrac12(\by-\bX\beta^{t})^{\top}(\by-\bX\beta^{t})\big)$.
\end{enumerate}

\subsection{Estimación de la media de $Y$ cuando $X=10$}
Sea $\mu_{10}=\E(Y\mid X=10)=\beta_0+\beta_1\cdot 10$. Necesitamos $p(\mu_{10}\mid\by)$, que se simula con Monte Carlo:
\[
\mu_{10}^{t}=\beta_0^{t}+\beta_1^{t}\times 10,\ t=1,\dots,10000\ \ (\text{muestras de }p(\mu_{10}\mid\by)).
\]

\subsection{Distribución predictiva a posteriori}
Sea $y^p$ el valor predictivo:
\[
p(y^p\mid\by)=\int p(y^p,\theta\mid\by)\,\dpar\theta=\int p(y^p\mid\theta,\by)\,p(\theta\mid\by)\,\dpar\theta.
\]
Si $y^p$ e $\by$ son independientes condicional a $\theta$, se reduce a
\[
p(y^p\mid\by)\propto\int p(y^p\mid\theta)\,p(\theta\mid\by)\,\dpar\theta.
\]
\textbf{Monte Carlo:} $\theta^t\sim p(\theta\mid\by)$ (Gibbs); $y_p^t\sim p(y^p\mid\theta=\theta^t)$.\\
\textbf{Ejemplo ($X=10$):} $Y_{10}\mid\beta,\sigma^2\sim\N(\beta_0+\beta_1\times 10,\sigma^2)$.
\begin{enumerate}[leftmargin=2em]
  \item $\beta^t,\sigma_t^2\sim p(\beta,\sigma^2\mid\by)$.
  \item $Y_{10}^t\sim\N(\beta_0^t+\beta_1^t\times 10,\ \sigma_t^2)$.
\end{enumerate}

% ============================================================
\section{Criterios para comparación de modelos: DIC}
% ============================================================
\textbf{DIC (Deviance Information Criterion).} Deviance: $D(\theta)=-2\log p(\by\mid\theta)$. Con $\hat\theta=\E(\theta\mid\by)$, $\hat D=D(\hat\theta)$, $\bar D=\E(D(\theta))$. Penalidad (número efectivo de parámetros):
\[
p_D=\bar D-\hat D,\qquad
\text{DIC}=\underbrace{\hat D}_{\text{ajuste}}+\underbrace{2p_D}_{\text{penaliza por complejidad}}.
\]
Con $\theta=(\beta,\sigma^2)$ y $y_i\overset{i.i.d.}{\sim}\N(X_i^{\top}\beta,\sigma^2)$:
\[
D(\theta)=-2\log p(\by\mid\beta,\sigma^2)=-2\log\prod_{i=1}^n\phi(y_i\mid X_i^{\top}\beta,\sigma^2)=-2\sum_{i=1}^n\log\phi(y_i\mid X_i^{\top}\beta,\sigma^2).
\]

\paragraph{Propiedad ($t$-Student).} Si $Z\sim\N(0,1)$ y $W\sim\Ga(\tfrac{\nu}{2},\tfrac{\nu}{2})$ independientes, entonces $T=\dfrac{Z}{\sqrt{W}}\sim t(0,1,\nu)$.

\begin{center}
\begin{tabular}{lc}
\hline
Modelo & DIC \\ \hline
Normal & $122.85$ \\
$t$-Student & $114.01$ \\ \hline
\end{tabular}
\end{center}

% ============================================================
\section{Modelo de regresión con errores $t$-Student}
% ============================================================
\subsection{Distribución $t$-Student}
Existe porque la normal asume que todos los datos son bien portados (sin valores extremos). Con outliers, la normal distorsiona la estimación; la $t$-Student tiene colas más pesadas. $Y\sim t(\mu,\sigma,\nu)$:
\begin{itemize}[leftmargin=2em]
  \item $\mu$: centro de la distribución.
  \item $\sigma$: escala.
  \item $\nu$: grados de libertad que controlan el grosor de las colas.
\end{itemize}
Cuando $\nu\to\infty$ la $t$-Student converge a una Normal; si $\nu$ es chico tiene colas pesadas.
\[
f_Y(y)=\frac{\Gamma((\nu+1)/2)}{\Gamma(\nu/2)\sqrt{\nu\pi}}\frac{1}{\sigma}\left(1+\frac{1}{\nu}\Big(\frac{y-\mu}{\sigma}\Big)^2\right)^{-\frac{\nu+1}{2}},\ y\in\R,\ \mu\in\R,\ \sigma>0,\ \nu>0,
\]
\[
\E[Y]=\mu\ (\nu>1),\qquad \Var[Y]=\frac{\nu}{\nu-2}\sigma^2\ (\nu>2).
\]

\paragraph{Estandarización.} Si $Z\sim t(0,1,\nu)$, entonces $Y=\mu+\sigma Z\sim t(\mu,\sigma,\nu)$, y $f_Y(y)=\tfrac{1}{\sigma}f_Z\big(\tfrac{y-\mu}{\sigma},\nu\big)$.

\paragraph{Representación estocástica.} Sean $Z\sim\N(0,1)$ y $W\sim\Ga(\nu/2,\nu/2)$ independientes; entonces $T=\dfrac{Z}{\sqrt{W}}\sim t(0,1,\nu)$.

\subsection{Modelo de regresión con errores $t$-Student}
\[
y_i=X_i^{\top}\beta+t_i,\qquad t_i\overset{ind}{\sim}t(0,1,\sigma^2).
\]
Usando la representación estocástica $t_i=\dfrac{\varepsilon_i}{\sqrt{w_i}}$ con $\varepsilon_i\sim\N(0,\sigma^2)$ y $w_i\sim\Ga(\nu/2,\nu/2)$ independientes, considerando $w_i$ como variable latente:
\[
y_i\mid w_i,\beta,\sigma^2\sim\N\!\left(X_i^{\top}\beta,\ \frac{\sigma^2}{w_i}\right),\qquad w_i\mid\beta,\sigma^2\sim\Ga(\nu/2,\nu/2).
\]
Para $i=1,\dots,n$, la verosimilitud de $\by=(y_1,\dots,y_n)$ y $\bW=(w_1,\dots,w_n)$ (regla del producto) es
\[
p(\by,\bW\mid\beta,\sigma^2)=\prod_{i=1}^n p(y_i\mid w_i,\beta,\sigma^2)\,p(w_i\mid\beta,\sigma^2)=\prod_{i=1}^n p(y_i\mid w_i,\beta,\sigma^2)\prod_{i=1}^n p(w_i\mid\beta,\sigma^2).
\]

\subsection{Priors y posterior}
Con $\nu$ fijo pequeño y prior no informativa $p(\beta,\sigma^2)\propto\dfrac{1}{\sigma^2}$:
\[
p(\beta,\sigma^2\mid\by)\propto\prod_{i=1}^n p(y_i\mid w_i,\beta,\sigma^2)\prod_{i=1}^n p(w_i\mid\beta,\sigma^2)\times\frac{1}{\sigma^2},
\]
donde
\[
p(y_i\mid w_i,\beta,\sigma^2)=\frac{w_i^{1/2}}{\sqrt{2\pi\sigma^2}}\exp\!\left(-\frac{w_i}{2\sigma^2}(y_i-X_i^{\top}\beta)^2\right),\qquad
p(w_i)=\frac{(\nu/2)^{\nu/2}}{\Gamma(\nu/2)}w_i^{\nu/2-1}e^{-(\nu/2)w_i}.
\]

\subsection{Condicionales completas (Gibbs)}
\paragraph{$\beta$.}
\[
p(\beta\mid\sigma^2,\omega,\by)\propto\prod_{i=1}^n\frac{w_i^{1/2}}{\sigma}\exp\!\left(-\frac{w_i}{2\sigma^2}(y_i-X_i^{\top}\beta)^2\right)\propto\exp\!\left(-\frac{1}{2\sigma^2}\sum_{i=1}^n w_i(y_i-X_i^{\top}\beta)^2\right).
\]
Con $\sum w_i(y_i-X_i^{\top}\beta)^2=(\by-\bX\beta)^{\top}\bW(\by-\bX\beta)$, expandiendo,
\[
\propto\exp\!\left(-\tfrac{1}{2\sigma^2}\big[\beta^{\top}\bX^{\top}\bW\bX\beta-2\beta^{\top}\bX^{\top}\bW\by\big]\right).
\]
Identificando: $\Sigma^{-1}=\dfrac{\bX^{\top}\bW\bX}{\sigma^2}\Rightarrow\Sigma=\sigma^2(\bX^{\top}\bW\bX)^{-1}$ y $\mu=(\bX^{\top}\bW\bX)^{-1}\bX^{\top}\bW\by$, entonces
\[
\boxed{\ \beta\mid\sigma^2,\omega,\by\sim\N\!\left((\bX^{\top}\bW\bX)^{-1}\bX^{\top}\bW\by,\ \sigma^2(\bX^{\top}\bW\bX)^{-1}\right)\ }
\]

\paragraph{$\sigma^2$.} Con $\text{SSE}_\omega=\sum_{i=1}^n w_i(y_i-X_i^{\top}\beta)^2$:
\[
p(\sigma^2\mid\beta,\omega,\by)\propto\left(\frac{1}{\sigma^2}\right)^{n/2+1}\exp\!\left(-\frac{\text{SSE}_\omega}{2\sigma^2}\right)
\ \Rightarrow\ \sigma^2\mid\beta,\omega,\by\sim\IG\!\left(\frac{n}{2},\ \frac{1}{2}\sum_{i=1}^n w_i(y_i-X_i^{\top}\beta)^2\right).
\]

\paragraph{$w_i$ (cada $i$ por separado).}
\begin{align*}
p(w_i\mid\beta,\sigma^2,y_i)&\propto p(y_i\mid w_i,\beta,\sigma^2)\,p(w_i)\\
&\propto\frac{w_i^{1/2}}{\sigma}\exp\!\left(-\frac{w_i}{2\sigma^2}(y_i-X_i^{\top}\beta)^2\right)\times w_i^{\nu/2-1}\exp\!\left(-\frac{\nu}{2}w_i\right)\\
&\propto w_i^{(\nu+1)/2-1}\exp\!\left(-w_i\left[\frac{(y_i-X_i^{\top}\beta)^2}{2\sigma^2}+\frac{\nu}{2}\right]\right),
\end{align*}
por lo tanto
\[
w_i\mid\beta,\sigma^2,y_i\sim\Ga\!\left(\frac{\nu+1}{2},\ \frac{\nu}{2}+\frac{(y_i-X_i^{\top}\beta)^2}{2\sigma^2}\right).
\]

\subsection{Algoritmo de Gibbs para el modelo $t$}
\begin{enumerate}[leftmargin=2em]
  \item Generar $w_i^{s+1}\sim\Ga\!\left(\tfrac{\nu+1}{2},\ \tfrac{\nu}{2}+\tfrac{(y_i-X_i^{\top}\beta^s)^2}{2\sigma_s^2}\right)$.
  \item Generar $\beta^{s+1}\sim\N\!\left((\bX^{\top}\bW^{s+1}\bX)^{-1}\bX^{\top}\bW^{s+1}\by,\ \sigma_s^2(\bX^{\top}\bW^{s+1}\bX)^{-1}\right)$.
  \item Generar $\sigma^2_{s+1}\sim\IG\!\left(\tfrac{n}{2},\ \tfrac12\sum_{i=1}^n w_i^{s+1}(y_i-X_i^{\top}\beta^{s+1})^2\right)$.
\end{enumerate}

% ============================================================
\section{Algoritmo de Metropolis-Hastings}
% ============================================================
$\theta^t=$ valor actual de la cadena. $q(\theta\mid\theta^t)$: distribución generadora de candidatos, condicional al valor actual.
\begin{enumerate}[leftmargin=2em]
  \item Generar un candidato $\theta^{\text{cand}}\sim q(\theta\mid\theta^t)$.
  \item Calcular la probabilidad de aceptación
  \[
  \alpha(\theta^{\text{cand}},\theta^t)=\min\left\{1,\ \frac{f(\theta^{\text{cand}})}{f(\theta^t)}\times\frac{q(\theta^t\mid\theta^{\text{cand}})}{q(\theta^{\text{cand}}\mid\theta^t)}\right\}.
  \]
  \item Aceptar el candidato con probabilidad $\alpha$: $\theta^{t+1}=\theta^{\text{cand}}$; rechazar con probabilidad $1-\alpha$: $\theta^{t+1}=\theta^t$.
\end{enumerate}
La secuencia $\theta^0,\theta^1,\theta^2,\dots$ converge a $f(\theta)$. Solo se requiere simular $q(\cdot\mid\cdot)$, que se puede escoger libremente. Aunque esté garantizada la convergencia, la elección de $q(\cdot\mid\cdot)$ es crucial para que el algoritmo sea eficiente.

\subsection{Metropolis-Hastings — paseo aleatorio}
Distribución generadora de candidatos: normal centrada en el valor actual, $\theta^{\text{cand}}\sim\N(\theta^t,\Sigma)$.
\begin{align*}
p(\theta^{\text{cand}}\mid\theta^t)&=\frac{1}{(2\pi)^{n/2}|\Sigma|^{1/2}}\exp\!\left(-\tfrac12(\theta^{\text{cand}}-\theta^t)^{\top}\Sigma^{-1}(\theta^{\text{cand}}-\theta^t)\right),\\
p(\theta^t\mid\theta^{\text{cand}})&=\frac{1}{(2\pi)^{n/2}|\Sigma|^{1/2}}\exp\!\left(-\tfrac12(\theta^t-\theta^{\text{cand}})^{\top}\Sigma^{-1}(\theta^t-\theta^{\text{cand}})\right).
\end{align*}
Notemos que $p(\theta^{\text{cand}}\mid\theta^t)=p(\theta^t\mid\theta^{\text{cand}})$: la propuesta es \textbf{simétrica} y se cancela en $\alpha$.
\begin{enumerate}[leftmargin=2em]
  \item $\theta^{\text{cand}}\sim\N(\theta^t,\Sigma)$.
  \item Generar $U\sim\text{U}(0,1)$.
  \item Si $U<\min\left\{1,\ \dfrac{f(\theta^{\text{cand}})}{f(\theta^t)}\right\}\Rightarrow\theta^{t+1}=\theta^{\text{cand}}$; en caso contrario $\theta^{t+1}=\theta^t$.
\end{enumerate}
En escala logarítmica (más estable numéricamente):
\[
\log(U)\le\log\!\left(\frac{f(\theta^{\text{cand}})}{f(\theta^t)}\right)\le\log f(\theta^{\text{cand}})-\log f(\theta^t).
\]
El algoritmo funciona bien si la tasa de aceptación no es muy alta (35\% a 40\%). $\Sigma$ debe seleccionarse apropiadamente.

\paragraph{Sugerencia de Gelman (1995).}
\[
\Sigma=c^2\bV,
\]
donde $\bV$ es la covarianza de la posterior y $c=\dfrac{2.4}{\sqrt{K}}$, con $K=$ número de parámetros. $\bV$ se puede aproximar con la negativa de la Hessiana de la log-posterior evaluada en el máximo posterior (matriz de información de Fisher).

\subsection{Ejemplo: Regresión Probit}
$y_i\overset{ind}{\sim}\text{Bernoulli}(\pi_i)$, $\pi_i=\Phi(X_i^{\top}\beta)$.
\[
p(\by\mid\beta)=\prod_{i=1}^n\pi_i^{y_i}(1-\pi_i)^{1-y_i}=\prod_{i=1}^n\Phi(X_i^{\top}\beta)^{y_i}\big[1-\Phi(X_i^{\top}\beta)\big]^{1-y_i}.
\]
Con $p(\beta)\propto 1$,
\[
p(\beta\mid\by)\propto p(\by\mid\beta)\propto\prod_{i=1}^n\Phi(X_i^{\top}\beta)^{y_i}\big[1-\Phi(X_i^{\top}\beta)\big]^{1-y_i},
\]
\[
\log p(\beta\mid\by)\propto\sum_{i=1}^n y_i\log\Phi(X_i^{\top}\beta)+(1-y_i)\log\big[1-\Phi(X_i^{\top}\beta)\big].
\]
En R: \verb|sum(dbinom(y, 1, pnorm(X %*% beta), log = TRUE))|.

% ============================================================
\section{Regresión Probit con variables latentes (Albert y Chib)}
% ============================================================
\subsection{Distribución normal truncada}
Sea $Z\sim\N(\mu,\sigma^2)\ind(a,b)$ truncada en $(a,b)$:
\[
f_Z(z)=\frac{\tfrac{1}{\sigma}\phi\!\big(\tfrac{z-\mu}{\sigma}\big)}{\Phi\!\big(\tfrac{b-\mu}{\sigma}\big)-\Phi\!\big(\tfrac{a-\mu}{\sigma}\big)},\quad a<z<b,
\]
donde el denominador es la constante normalizadora que garantiza que la densidad integre $1$; $\phi(\cdot)$ y $\Phi(\cdot)$ son la densidad y CDF normal estándar. Notemos que $f_Z(z)\propto\tfrac{1}{\sigma}\phi\!\big(\tfrac{z-\mu}{\sigma}\big)\ind(a<z<b)$.

\subsection{Modelo Probit con latentes}
$y_i\sim\text{Bernoulli}(\pi_i)$, $\pi_i=\Phi(X_i^{\top}\beta)$ (CDF de la normal estándar mapea a $(0,1)$). La verosimilitud de $\by=(y_1,\dots,y_n)^{\top}$ es
\[
p(\by\mid\beta)=\prod_{i=1}^n\big[\Phi(X_i^{\top}\beta)\big]^{y_i}\big[1-\Phi(X_i^{\top}\beta)\big]^{1-y_i}.
\]
Con prior no informativa $p(\beta)\propto 1$, la posterior $p(\beta\mid\by)\propto\prod_{i=1}^n[\Phi(X_i^{\top}\beta)]^{y_i}[1-\Phi(X_i^{\top}\beta)]^{1-y_i}$ no tiene forma conocida. Se usa la propuesta de Albert y Chib (1993), que introduce variables latentes que vinculan la variable respuesta binaria con un modelo lineal subyacente:
\[
Z_i=X_i^{\top}\beta+\varepsilon_i,\ \varepsilon_i\sim\N(0,1),\qquad Y_i=\begin{cases}1 & \text{si }Z_i>0\\ 0 & \text{si }Z_i\le0\end{cases}
\]
Notemos que
\[
\pi_i=P(Y_i=1)=P(Z_i>0)=P(\varepsilon_i>-X_i^{\top}\beta)=P(\varepsilon_i<X_i^{\top}\beta)=\Phi(X_i^{\top}\beta),
\]
por simetría de la normal estándar. Con $\sigma=1$, $p(z_i\mid\beta)=\phi(z_i-X_i^{\top}\beta)$. La verosimilitud aumentada es
\[
p(\by,\bz\mid\beta)=\prod_{i=1}^n p(z_i\mid\beta)\,p(y_i\mid z_i,\beta)
\propto\prod_{i=1}^n\phi(z_i-X_i^{\top}\beta)\big[\ind(z_i>0)\ind(y_i=1)+\ind(z_i\le0)\ind(y_i=0)\big].
\]

\subsection{Condicionales (Gibbs)}
\paragraph{$\beta$.}
\begin{align*}
p(\beta\mid\bz,\by)&\propto\prod_{i=1}^n\phi(z_i-X_i^{\top}\beta)
\propto\exp\!\left(-\tfrac12\sum_{i=1}^n(z_i-X_i^{\top}\beta)^2\right)
\propto\exp\!\left(-\tfrac12(\bz-\bX\beta)^{\top}(\bz-\bX\beta)\right)\\
&\propto\exp\!\left(-\tfrac12\big[\beta^{\top}\bX^{\top}\bX\beta-2\beta^{\top}\bX^{\top}\bz\big]\right),
\end{align*}
por lo tanto $\beta\mid\bz,\by\sim\N\big((\bX^{\top}\bX)^{-1}\bX^{\top}\bz,\ (\bX^{\top}\bX)^{-1}\big)$.

\paragraph{$z_i$.} Si $y_i=1$:
\begin{multline*}
p(z_i\mid\beta,y_i=1)\propto p(z_i\mid\beta)P(y_i=1\mid z_i,\beta)\propto\phi(z_i-X_i^{\top}\beta)\ind(z_i>0)\\
\Rightarrow\ z_i\mid\beta,y_i=1\sim\N(X_i^{\top}\beta,1)\ind(0,\infty).
\end{multline*}
Si $y_i=0$:
\[
p(z_i\mid\beta,y_i=0)\propto\phi(z_i-X_i^{\top}\beta)\ind(z_i\le0)\ \Rightarrow\ z_i\mid\beta,y_i=0\sim\N(X_i^{\top}\beta,1)\ind(-\infty,0).
\]

\subsection{Algoritmo de Gibbs}
\begin{enumerate}[leftmargin=2em]
  \item Generar $\begin{cases}z_i^{s+1}\sim\N(X_i^{\top}\beta^s,1)\ind(0,\infty) & \text{si }y_i=1\\ z_i^{s+1}\sim\N(X_i^{\top}\beta^s,1)\ind(-\infty,0) & \text{si }y_i=0\end{cases}$
  \item Generar $\beta^{s+1}\sim\N\big((\bX^{\top}\bX)^{-1}\bX^{\top}\bz^{s+1},\ (\bX^{\top}\bX)^{-1}\big)$.
\end{enumerate}

% ============================================================
\section{Modelo para datos con censura intervalar (brecha salarial)}
% ============================================================
\textbf{Contexto:} la banda salarial se postula como datos censurados. Sea $Y_i$ el número de intervalo de ingreso observado del $i$-ésimo individuo y $Z_i$ el ingreso no observado:
\[
\begin{aligned}
Y_i=1 &\iff 850\le Z_i\le 999, & Y_i=2 &\iff 1000\le Z_i\le 2000,\\
Y_i=3 &\iff 2001\le Z_i\le 3000, & Y_i=4 &\iff 3001\le Z_i\le 4000,\\
Y_i=5 &\iff 4001\le Z_i\le 5000, & Y_i=6 &\iff 5001\le Z_i\le\infty.
\end{aligned}
\]
Se propone que el ingreso no observado se modele como log-normal (todos sus valores son positivos): $Z_i\overset{ind}{\sim}\text{LN}(\mu_i,\sigma^2)$, de manera que
\[
\text{Me}(Z_i)=e^{\mu_i},\quad \E(Z_i)=e^{\mu_i+\sigma^2/2},\quad \Var(Z_i)=e^{2\mu_i+\sigma^2}(e^{\sigma^2}-1).
\]
Escribiendo el modelo en términos de $w_i=\log(Z_i)$: los intervalos $[l_j,u_j]$ ($j=1,\dots,6$) se transforman como $l_j=\log(\cdot)$, $u_j=\log(\cdot)$, con $w_i\overset{ind}{\sim}\N(\mu_i,\sigma^2)$ y $\mu_i=\beta_0+\beta_1\,\text{sexo}_i$.

\paragraph{Conjunta.}
\begin{align*}
p(\by,\omega\mid\beta,\sigma^2)&=\prod_{i=1}^n p(w_i\mid\beta,\sigma^2)\,p(y_i\mid w_i,\beta,\sigma^2)\\
&=\prod_{i=1}^n\frac{1}{\sqrt{2\pi}}\left(\frac{1}{\sigma^2}\right)^{1/2}\exp\!\left(-\frac{1}{2\sigma^2}(w_i-X_i^{\top}\beta)^2\right)\left[\sum_{j=1}^6\ind(l_j\le w_i\le u_j)\ind(y_i=j)\right].
\end{align*}
Con prior no informativa $\propto 1/\sigma^2$, la posterior es
\[
p(\beta,\sigma^2,\omega\mid\by)\propto\prod_{i=1}^n\frac{1}{\sqrt{2\pi}}\Big(\tfrac{1}{\sigma^2}\Big)^{1/2}\exp\!\left(-\tfrac{1}{2\sigma^2}(w_i-X_i^{\top}\beta)^2\right)\left[\sum_{j=1}^6\ind(l_j\le w_i\le u_j)\ind(y_i=j)\right]\times\frac{1}{\sigma^2}.
\]

\subsection{Condicionales (Gibbs)}
\paragraph{$\beta$.}
\[
p(\beta\mid\sigma^2,\omega,\by)\propto\exp\!\left(-\tfrac{1}{2\sigma^2}(\omega-\bX\beta)^{\top}(\omega-\bX\beta)\right)
\ \Rightarrow\ \beta\mid\bz,\by\sim\N\big((\bX^{\top}\bX)^{-1}\bX^{\top}\omega,\ \sigma^2(\bX^{\top}\bX)^{-1}\big).
\]

\paragraph{$\sigma^2$.} Con $\text{SSE}=\sum_{i=1}^n(w_i-X_i^{\top}\beta)^2$:
\[
\sigma^2\mid\beta,\omega,\by\sim\IG\!\left(\frac{n}{2},\ \frac{1}{2}\sum_{i=1}^n(w_i-X_i^{\top}\beta)^2\right).
\]

\paragraph{$w_i$ (latentes).}
\begin{align*}
p(w_i\mid\beta,\sigma^2,y_i)&\propto p(w_i\mid\beta,\sigma^2)\,p(y_i\mid w_i,\beta,\sigma^2)\\
&\propto\frac{1}{\sigma}\phi\!\left(\frac{w_i-X_i^{\top}\beta}{\sigma}\right)\ind(l_i\le w_i\le u_i),
\end{align*}
por lo tanto
\[
w_i\mid\beta,\sigma^2,y_i\sim\N(X_i^{\top}\beta,\sigma^2)\ind(l_i,u_i).
\]

\subsection{Algoritmo de Gibbs}
\begin{enumerate}[leftmargin=2em]
  \item Generar $w_i^{s+1}\sim\N(X_i^{\top}\beta^s,\sigma_s^2)\ind(l_i,u_i)$, $i=1,\dots,n$.
  \item Generar $\beta^{s+1}\sim\N\big((\bX^{\top}\bX)^{-1}\bX^{\top}\omega^{s+1},\ \sigma_s^2(\bX^{\top}\bX)^{-1}\big)$.
  \item Generar $\sigma^2_{s+1}\sim\IG\big(\tfrac{n}{2},\ \tfrac12\sum_{i=1}^n(w_i^{s+1}-X_i^{\top}\beta^{s+1})^2\big)$.
\end{enumerate}

\subsection{Estimación de la brecha salarial por género}
Como asumimos log-normal sobre el salario real $Z_i$:
\[
\frac{\E(\text{Salario}\mid\text{Mujer})}{\E(\text{Salario}\mid\text{Hombre})}=\frac{\E[Z_i\mid X_i=1]}{\E[Z_i\mid X_i=0]}=\frac{\exp(\beta_0+\beta_1+\sigma^2/2)}{\exp(\beta_0+\sigma^2/2)}=\exp(\beta_1).
\]
Parámetro de brecha salarial por género: $\delta=1-\exp(\beta_1)$. Se obtienen los momentos y los IC para hacer inferencia a partir del muestreo de los parámetros.

% ============================================================
\section{Modelo de regresión con punto de cambio (discreto)}
% ============================================================
Se desea estimar el valor de $\tau$ de la razón agua/cemento a partir del cual la resistencia a la compresión del concreto decae rápidamente a medida que aumenta dicha razón.
\[
Y_i\overset{ind}{\sim}\N(\mu_i,\sigma^2),\qquad \mu_i=\beta_0+\beta_1 X_i+\beta_2(X_i-\tau)\ind(X_i>\tau),
\]
con $i=1,\dots,n$, $\beta=(\beta_0,\beta_1,\beta_2)^{\top}$ los coeficientes de regresión, $\sigma^2>0$ la varianza de los errores y $\ind(\cdot)$ la función indicadora. El modelo genera dos regímenes lineales:
\[
\mu_i=\begin{cases}
\beta_0+\beta_1 X_i, & X_i\le\tau,\\[2pt]
(\beta_0-\beta_2\tau)+(\beta_1+\beta_2)X_i, & X_i>\tau,
\end{cases}
\]
es decir, un cambio de intercepto y de pendiente en el punto $\tau$. El punto de cambio $\tau$ es un parámetro desconocido, discreto con rango $\tau=50,51,\dots,81,82$.

\paragraph{Prior.} $\theta=(\beta,\sigma^2,\tau)^{\top}$, $p(\theta)=p(\beta)p(\sigma^2)p(\tau)$, con
\[
\beta\sim\N(\beta_0,\Sigma_0),\qquad \sigma^2\sim\IG(a_0,b_0),\qquad \tau\sim\text{Uniforme-Discreta}\{50,\dots,82\}.
\]

\paragraph{Verosimilitud.} $\by\mid\theta\sim\N(X_{(\tau)}\beta,\sigma^2\bI)$, con $X_{(\tau)}$ matriz con filas $X_i(\tau)=(1,X_i,(X_i-\tau)\ind(X_i>\tau))$:
\[
p(\by\mid\theta)=\left(\frac{1}{2\pi}\right)^{n/2}\left(\frac{1}{\sigma^2}\right)^{n/2}\exp\!\left(-\frac{1}{2\sigma^2}(\by-X_{(\tau)}\beta)^{\top}(\by-X_{(\tau)}\beta)\right).
\]
La posterior es
\begin{align*}
p(\theta\mid\by)&\propto p(\by\mid\theta)p(\beta)p(\sigma^2)p(\tau)\\
&\propto\Big(\tfrac{1}{\sigma^2}\Big)^{n/2}\exp\!\left(-\tfrac{1}{2\sigma^2}(\by-X_{(\tau)}\beta)^{\top}(\by-X_{(\tau)}\beta)\right)\exp\!\left(-\tfrac12(\beta-\beta_0)^{\top}\Sigma_0^{-1}(\beta-\beta_0)\right)\\
&\quad\times\Big(\tfrac{1}{\sigma^2}\Big)^{a_0+1}\exp\!\left(-\tfrac{b_0}{\sigma^2}\right).
\end{align*}

\subsection{Condicionales (Gibbs)}
\paragraph{$\beta$.}
\[
p(\beta\mid\sigma^2,\tau,\by)\propto\exp\!\left(-\tfrac{1}{2}\Big[\beta^{\top}\Big(\Sigma_0^{-1}+\tfrac{X_{(\tau)}^{\top}X_{(\tau)}}{\sigma^2}\Big)\beta-2\beta^{\top}\Big(\Sigma_0^{-1}\beta_0+\tfrac{X_{(\tau)}^{\top}\by}{\sigma^2}\Big)\Big]\right),
\]
por lo que
\[
\Sigma_\beta^{-1}=\Sigma_0^{-1}+\frac{X_{(\tau)}^{\top}X_{(\tau)}}{\sigma^2},\qquad
\mu_\beta=\Sigma_\beta\Big(\Sigma_0^{-1}\beta_0+\frac{X_{(\tau)}^{\top}\by}{\sigma^2}\Big),
\qquad
\beta\mid\sigma^2,\tau,\by\sim\N(\mu_\beta,\Sigma_\beta).
\]

\paragraph{$\sigma^2$.} Con $\text{SSE}=(\by-X_{(\tau)}\beta)^{\top}(\by-X_{(\tau)}\beta)$:
\[
\sigma^2\mid\beta,\tau,\by\sim\IG\!\left(a_0+\frac{n}{2},\ b_0+\frac{1}{2}\text{SSE}\right).
\]

\paragraph{$\tau$.}
\[
p(\tau\mid\beta,\sigma^2,\by)\propto p(\by\mid\beta,\sigma^2,\tau)\,p(\tau)\propto\Big(\tfrac{1}{\sigma^2}\Big)^{n/2}\exp\!\left(-\tfrac{1}{2\sigma^2}(\by-X_{(\tau)}\beta)^{\top}(\by-X_{(\tau)}\beta)\right),
\]
que se evalúa en cada valor discreto del rango de $\tau$ y se normaliza.

\paragraph{Algoritmo.}
\begin{enumerate}[leftmargin=2em]
  \item Generar $\beta^{s+1}\sim\N(\cdot,\cdot)$.
  \item Generar $\sigma^2_{s+1}\sim\IG(\cdot,\cdot)$.
  \item Generar $\tau^{s+1}$ de $p(\tau\mid\beta^{s+1},\sigma^{s+1},\by)$.
\end{enumerate}

% ============================================================
\section{Metropolis-Hastings con punto de cambio continuo}
% ============================================================
Mismo modelo con $\mu_i=\beta_0+\beta_1 X_i+\beta_2(X_i-\tau)\ind(X_i>\tau)$, pero ahora $\tau$ es continuo, $0<\tau<1$.

\paragraph{Prior.} $\theta=(\beta,\sigma^2,\tau)^{\top}$, $p(\theta)=p(\beta)p(\sigma^2)p(\tau)$, con
\[
\beta\sim\N(\beta_0,\Sigma_0),\qquad\sigma^2\sim\IG(a_0,b_0),\qquad\tau\sim\text{Beta}(c_0,d_0).
\]

\paragraph{Verosimilitud y posterior.} $\by\mid\theta\sim\N(X_{(\tau)}\beta,\sigma^2\bI)$, con $X_i(\tau)=(1,X_i,(X_i-\tau)\ind(X_i>\tau))$:
\begin{align*}
p(\theta\mid\by)&\propto\Big(\tfrac{1}{\sigma^2}\Big)^{n/2}\exp\!\left(-\tfrac{1}{2\sigma^2}(\by-X_{(\tau)}\beta)^{\top}(\by-X_{(\tau)}\beta)\right)\times\exp\!\left(-\tfrac12(\beta-\beta_0)^{\top}\Sigma_0^{-1}(\beta-\beta_0)\right)\\
&\quad\times\Big(\tfrac{1}{\sigma^2}\Big)^{a_0+1}\exp\!\left(-\tfrac{b_0}{\sigma^2}\right)\times\tau^{c_0-1}(1-\tau)^{d_0-1}.
\end{align*}

\subsection{Condicionales}
\paragraph{$\beta$.}
\[
\beta\mid\sigma^2,\tau,\by\sim\N\!\left(\Big(\Sigma_0^{-1}+\tfrac{X_{(\tau)}^{\top}X_{(\tau)}}{\sigma^2}\Big)^{-1}\Big(\Sigma_0^{-1}\beta_0+\tfrac{X_{(\tau)}^{\top}\by}{\sigma^2}\Big),\ \Big(\Sigma_0^{-1}+\tfrac{X_{(\tau)}^{\top}X_{(\tau)}}{\sigma^2}\Big)^{-1}\right).
\]

\paragraph{$\sigma^2$.} Con $\text{SSE}=(\by-X_{(\tau)}\beta)^{\top}(\by-X_{(\tau)}\beta)=\sum_{i=1}^n(y_i-X_i(\tau)^{\top}\beta)^2$:
\[
\sigma^2\mid\beta,\tau,\by\sim\IG\!\left(a_0+\frac{n}{2},\ b_0+\frac{1}{2}(\by-X_{(\tau)}\beta)^{\top}(\by-X_{(\tau)}\beta)\right).
\]

\paragraph{$\tau$.}
\[
p(\tau\mid\beta,\sigma^2,\by)\propto\prod_{i=1}^n\frac{1}{\sqrt{2\pi}\sigma}\exp\!\left(-\tfrac{1}{2\sigma^2}(y_i-X_i(\tau)^{\top}\beta)^2\right)\times\frac{\Gamma(c_0+d_0)}{\Gamma(c_0)\Gamma(d_0)}\tau^{c_0-1}(1-\tau)^{d_0-1},
\]
que no tiene forma cerrada. Para simular de ella se utiliza M-H con paseo aleatorio normal.

\paragraph{Construcción de datos en R.}
\begin{itemize}[leftmargin=2em]
  \item Variables continuas: \verb|X = cbind(1, X, (X - tau)*(X > tau))|.
  \item Variables categóricas con \verb|model.matrix|:\\
  \verb|d$rank = as.factor(d$rank)|\\
  \verb|X = model.matrix(~ gre + gpa + rank, d)|.
\end{itemize}

% ============================================================
\section{Distribuciones conjugadas a recordar}
% ============================================================
\paragraph{Poisson--Gamma.} $Y_i\mid\lambda\sim\text{Poisson}(\lambda)$, $\lambda\sim\Ga(a,b)$:
\[
p(\by\mid\lambda)=\prod_{i=1}^n\frac{\lambda^{y_i}e^{-\lambda}}{y_i!}\propto\lambda^{\sum y_i}e^{-n\lambda},
\qquad \lambda\mid\by\sim\Ga\!\left(a+\sum_{i=1}^n y_i,\ b+n\right),
\]
con $\E[Y]=\lambda$, $\Var[Y]=\lambda$.

\paragraph{Exponencial--Gamma.} $Y_i\mid\lambda\sim\text{Exp}(\lambda)$, $\lambda\sim\Ga(a,b)$:
\[
p(\by\mid\lambda)=\prod_{i=1}^n\lambda e^{-\lambda y_i}=\lambda^n e^{-\lambda\sum y_i},
\qquad \lambda\mid\by\sim\Ga\!\left(a+n,\ b+\sum y_i\right).
\]

\paragraph{Bernoulli--Beta.} $Y_i\mid p\sim\text{Bernoulli}(p)$, $p\sim\text{Beta}(a,b)$:
\[
p(\by\mid p)=\prod_{i=1}^n p^{y_i}(1-p)^{1-y_i}=p^{\sum y_i}(1-p)^{n-\sum y_i},
\qquad p\mid\by\sim\text{Beta}\!\left(a+\sum y_i,\ b+n-\sum y_i\right).
\]

\paragraph{Binomial--Beta.} $Y_i\mid p\sim\text{Binomial}(n_i,p)$, $p\sim\text{Beta}(a,b)$:
\[
p(\by\mid p)\propto p^{\sum y_i}(1-p)^{\sum(n_i-y_i)},
\qquad p\mid\by\sim\text{Beta}\!\left(a+\sum y_i,\ b+\sum(n_i-y_i)\right).
\]

% ============================================================
\section{Análisis de convergencia}
% ============================================================
La convergencia de una cadena de Markov hacia su distribución estacionaria (posterior) se evalúa con pruebas estadísticas o inspección visual de las cadenas.
\begin{itemize}[leftmargin=2em]
  \item \textbf{Gráfico de trayectorias:} al alcanzar convergencia, las cadenas se mezclan y siguen una trayectoria estable y densa (banda horizontal).
  \item \textbf{Gráficos de autocorrelación:} las autocorrelaciones deben tender a decaer al aumentar los rezagos. Autocorrelación fuerte no impide la convergencia, pero reduce la eficiencia.
\end{itemize}
Se puede hacer inferencia si el gráfico de trayectorias alcanza un estado estacionario y las autocorrelaciones decaen rápido. Si no hay convergencia, el primer paso es incrementar el burn-in y realizar más simulaciones. Si hay convergencia pero autocorrelaciones altas, se puede guardar solo una de cada $k$ iteraciones (\emph{thinning}, período $\texttt{thin}$); esto solo sirve para ahorrar memoria.

\subsection{Prueba de Geweke}
Prueba de diferencia de medias entre la parte inicial y final de la cadena. Se compara el 10\% inicial frente al 50\% final, descartando iteraciones intermedias.
\[
Z=\frac{\bar\theta_A-\bar\theta_B}{\sqrt{\tfrac{\hat S_A(0)}{N_A}+\tfrac{\hat S_B(0)}{N_B}}}\sim\N(0,1),
\]
donde $\bar\theta_A,\bar\theta_B$ son las medias de las partes inicial y final, $N_A,N_B$ el número de iteraciones, y $\hat S_A(0)$ y $\hat S_B(0)$ son las densidades espectrales en frecuencia cero (varianzas asintóticas ajustadas por autocorrelación):
\[
S(0)=\sigma^2\left(1+2\sum_{k=1}^\infty\rho_k\right).
\]
Si $|Z|>1.96$ (nivel 5\%) se rechaza la hipótesis nula: la cadena aún no ha convergido.

\subsection{Prueba de Gelman y Rubin (PSRF)}
Compara simultáneamente múltiples cadenas corridas en paralelo. Con $m$ cadenas de $S$ simulaciones:
\[
W=\frac{1}{m}\sum_{j=1}^m\left[\frac{1}{S-1}\sum_{s}(\theta_s^j-\bar\theta^j)^2\right]\ \ (\text{varianza dentro de las cadenas}),
\]
\[
B=\frac{S}{m-1}\sum_{j=1}^m(\bar\theta^j-\bar\theta)^2\ \ (\text{varianza entre cadenas}),
\]
donde $\bar\theta$ es la media general. La estadística es
\[
\hat R=\sqrt{\frac{\tfrac{S-1}{S}W+\tfrac{1}{S}B}{W}}.
\]
Se espera $\hat R\approx 1$; en la práctica $\hat R<1.05$ indica convergencia.

\subsection{Tamaño efectivo de muestra}
\[
N_{\text{eff}}=\frac{N}{\sum_{t=-\infty}^{\infty}\rho_t}=\frac{N}{1+2\sum_{t=1}^{\infty}\rho_t},
\]
donde $\rho_t$ es la autocorrelación en el lag $t$ y $N$ el número de simulaciones. A mayor autocorrelación, menor $N_{\text{eff}}$.

% ============================================================
\section{Slice Sampling}
% ============================================================
Quiero generar muestras de $p(\theta)$, pero tiene forma rara y no puedo aplicar Gibbs directo. Idea geométrica: bajo la curva de $p(\theta)$ hay una región
\[
\mathcal{R}=\{(\theta,u):0\le u\le p(\theta)\}.
\]
Generar puntos uniformes debajo de la curva es equivalente a muestrear de $p(\theta)$. El problema es que generar puntos uniformes bajo la curva no es fácil (habría que conocer el área en cada región). En cada iteración, con punto actual $\theta^s$:
\begin{enumerate}[leftmargin=2em]
  \item Dado $\theta^s$, la altura de la curva es $p(\theta^s)$. Se tira $u^{s+1}\sim\text{U}(0,p(\theta^s))$: define un nivel de corte horizontal (la región válida).
  \item Nos movemos horizontalmente dentro del slice: $\theta^{s+1}\sim\text{U}\big(\{\theta:p(\theta)\ge u^{s+1}\}\big)$.
\end{enumerate}
Para encontrar los bordes del slice hay que resolver $p(\theta)=u^{s+1}$, que puede ser imposible analíticamente. Neal (2003) propone un mecanismo iterativo en dos fases:
\begin{itemize}[leftmargin=2em]
  \item \textbf{Stepping out (expansión):} parte de un intervalo pequeño alrededor de $\theta^s$ y lo expande hacia izquierda y derecha en pasos de tamaño $\omega$ hasta que ambos extremos caigan fuera del slice ($p<u^{s+1}$).
  \item \textbf{Shrinkage:} toma un candidato $\theta'$ uniforme del intervalo expandido. Si $p(\theta')<u^{s+1}$, contrae el intervalo hacia $\theta^s$ por ese lado y reintenta. Cada rechazo achica el intervalo, garantizando que eventualmente caiga dentro.
\end{itemize}

\subsection{Automated Factor Slice Sampling (AFSS)}
El algoritmo anterior opera una dimensión a la vez y funciona cuando los parámetros son independientes. Cuando dos parámetros están fuertemente correlacionados, la posterior tiene forma de elipse diagonal y si muestreamos en dirección horizontal o vertical el slice queda angostísimo. La solución es rotar el sistema de referencia para que las direcciones de exploración estén alineadas con el elipse: se realiza slice sampling con los autovectores $\Gamma_k$ de la matriz de covarianza $\Lambda$ de la posterior. $\theta^*+\eta_k\Gamma_k$ mueve $\theta^*$ en la dirección $\Gamma_k$ una distancia $\eta_k$.
\begin{enumerate}[leftmargin=2em]
  \item Simular $u^{s+1}\sim\text{U}(0,p(\theta^s))$.
  \item Inicializar $\theta^*=\theta^s$. Para cada autovector $\Gamma_k$: generar $\eta_k^*\sim\text{U}\big(\{\eta_k:p(\theta^*+\eta_k\Gamma_k)\ge u^{s+1}\}\big)$ (con shrinkage y stepping out) y actualizar $\theta^*=\theta^*+\eta_k^*\Gamma_k$.
  \item $\theta^{s+1}=\theta^*$.
\end{enumerate}
$\Lambda$ se estima directamente con las muestras del burn-in. \textbf{Limitaciones:} AFSS asume dependencia lineal entre parámetros (si es no lineal, el slice sigue angosto); estimar y descomponer $\Lambda$ de tamaño $K\times K$ tiene costo computacional alto, por lo que se aplica a bloques pequeños de parámetros a priori correlacionados.

\subsection{Barker Proposal}
Es un método de Metropolis-Hastings: propone un candidato y lo acepta o rechaza. La diferencia clave con M-H estándar (propuesta simétrica): en Barker se usa el gradiente de $\log p(\theta)$ para sesgar la propuesta hacia regiones de alta densidad.
\[
z_k\sim\N(0,\sigma^2)\ (\text{negativo o positivo}),\qquad
q(\theta_k^s,z_k)=\frac{1}{1+\exp(-z_k\,\partial_k\log p(\theta^s))}\in(0,1).
\]
Si el gradiente es positivo (la densidad crece hacia la derecha) y $z_k$ es positivo, la sigmoide se acerca a $1$. La dirección $b_k\in\{-1,1\}$ se asigna $b_k=1$ con probabilidad $q(\theta_k^s,z_k)$ y $b_k=-1$ con probabilidad $1-q(\theta_k^s,z_k)$. El candidato es $\theta_k^{\text{cand}}=\theta_k^s+b_k z_k$. Se acepta con
\[
\alpha^B(\theta^s,\theta^{\text{cand}})=\min\!\left(1,\ \frac{p(\theta^{\text{cand}})}{p(\theta^s)}\times\prod_{k=1}^K\frac{1+\exp\big((\theta_k^s-\theta_k^{\text{cand}})\partial_k\log p(\theta^s)\big)}{1+\exp\big((\theta_k^{\text{cand}}-\theta_k^s)\partial_k\log p(\theta^{\text{cand}})\big)}\right).
\]
El ratio no es $1$ como en RW porque la propuesta está sesgada.

% ============================================================
\section{Estructura Nimble (ejemplo: Betting / dosis-respuesta)}
% ============================================================
\begin{verbatim}
# Datos: insectos muertos (v. aleatorias del modelo)
data      <- list(Y = c(6, 13, 18, 28, 52, 53, 61, 60))
constants <- list(
  N = 8,                      # numero de experimentos
  X = c(...),                 # dosis
  n = c(59, 60, ...)          # insectos expuestos (valores fijos)
)

# Modelo
code <- nimbleCode({
  for (i in 1:N) {                        # verosimilitud por experimento
    Y[i]      ~ dbinom(size = n[i], prob = p[i])   # muertos ~ Binomial
    logit(p[i]) <- alpha + beta * X[i]             # dosis -> p[i]
  }
  alpha ~ dnorm(0, sd = 100)              # priors
  beta  ~ dnorm(0, sd = 100)
})

# Valores iniciales
inits <- function() list(alpha = 0, beta = 0)

# Construir y compilar
model   <- nimbleModel(code, constants, data, inits())
C_model <- compileNimble(model)

# Configurar MCMC
mcmcConf <- configureMCMC(model, monitors = c("alpha", "beta"))
mcmc     <- buildMCMC(mcmcConf)
C_mcmc   <- compileNimble(mcmc, project = model)

# Correr la simulacion
samples <- runMCMC(C_mcmc,
  niter        = 10000,   # iteraciones
  nburnin      = 2000,    # burn-in
  nchains      = 4,       # cadenas independientes
  summary      = TRUE,    # media, desviacion y cuantiles
  WAIC         = TRUE,    # metrica de comparacion de modelos
  samplesAsCodaMCMC = TRUE)

# Diagnosticos
mcmc_traces(samples$samples)   # trayectorias
mcmc_acf(samples$samples)      # autocorrelaciones
\end{verbatim}

\end{document}
